% ============================================================
% DATA ANALYSIS WITH PYTHON FOR HEALTH SPECIALISTS
% Lecture Notes (30 hours) — English version
% Compilation: xelatex -shell-escape notes.tex (2 passes)
% ============================================================
\documentclass[12pt,a4paper,twoside]{book}

% --- Encoding and language (XeLaTeX) ---
\usepackage{fontspec}
\usepackage{polyglossia}
\setmainlanguage{english}

% --- Page layout ---
\usepackage[margin=2.5cm]{geometry}
\usepackage{fancyhdr}
\pagestyle{fancy}
\fancyhead[LE]{\leftmark}
\fancyhead[RO]{\rightmark}
\fancyfoot[C]{\thepage}
\setlength{\headheight}{14.5pt}

% --- Mathematics ---
\usepackage{amsmath,amssymb,amsthm,mathtools}
\usepackage{bm}

% --- Graphics ---
\usepackage{tikz,pgfplots}
\pgfplotsset{compat=1.18}
\usetikzlibrary{patterns,arrows.meta,calc,positioning,shapes,fit,backgrounds}

% --- Code ---
\usepackage{minted}
\setminted{fontsize=\small,frame=lines,framesep=2mm,breaklines=true}
\setminted[python]{python3=true}

% --- Tables ---
\usepackage{booktabs,array,multirow,longtable}

% --- Colors and boxes ---
\usepackage[dvipsnames,svgnames]{xcolor}
\definecolor{steelblue}{RGB}{70,130,180}
\definecolor{healthgreen}{RGB}{34,139,34}
\definecolor{warnorange}{RGB}{255,140,0}
\usepackage{tcolorbox}
\tcbuselibrary{theorems,breakable,skins}

% --- Custom environments ---
\newtcolorbox{clinicalnote}{
  colback=healthgreen!5,colframe=healthgreen!70,
  title={\faStethoscope\ Clinical context},
  fonttitle=\bfseries,breakable
}
\newtcolorbox{pythontip}{
  colback=steelblue!5,colframe=steelblue!70,
  title={\faPython\ Python tip},
  fonttitle=\bfseries,breakable
}
\newtcolorbox{warningbox}{
  colback=warnorange!5,colframe=warnorange!70,
  title={\faExclamationTriangle\ Caution},
  fonttitle=\bfseries,breakable
}
\newtcolorbox{exercise}{
  colback=gray!5,colframe=gray!50,
  title={Exercise},
  fonttitle=\bfseries,breakable
}

% --- Links ---
\usepackage{hyperref}
\hypersetup{colorlinks=true,linkcolor=NavyBlue,citecolor=ForestGreen,urlcolor=Maroon}

% --- Bibliography and index ---
\usepackage[numbers]{natbib}
\usepackage{makeidx}
\makeindex

% --- Icons ---
\usepackage{fontawesome5}

% --- Theorem environments ---
\newtheorem{definition}{Definition}[chapter]
\newtheorem{example}{Example}[chapter]

% ============================================================
\title{
  \vspace{-2cm}
  {\Huge\bfseries Data Analysis with Python}\\[6pt]
  {\LARGE for Health Specialists}\\[12pt]
  {\large A 30-hour practical course}\\[24pt]
  {\normalsize Ya\'{e} Ulrich Gaba}\\[4pt]
  {\small AIRINA Labs}
}
\date{2026}
\author{}

% ============================================================
\begin{document}

% --- Title page ---
\frontmatter
\maketitle

% --- Table of contents ---
\tableofcontents
\newpage

% --- Preface ---
\chapter*{Preface}
\addcontentsline{toc}{chapter}{Preface}

This course is designed for health professionals --- physicians, nurses, epidemiologists, public health researchers, and health policy analysts --- who need to analyze data but have little or no programming experience. You do not need to become a software developer. You need to become someone who can load a dataset, clean it, explore it, test a hypothesis, build a simple predictive model, and communicate the results clearly.

We use Python because it is free, widely adopted in health research, and has an ecosystem of libraries specifically designed for the kinds of analyses health professionals need: survival analysis, epidemiological measures, geospatial mapping, and clinical prediction.

Every chapter follows the same rhythm: a brief explanation of the concept, a worked example on real health data, and exercises you complete in Jupyter notebooks. All datasets used in this course are freely available from the WHO, CDC, UCI Machine Learning Repository, and other open sources.

\paragraph{What this course is not.} This is not a statistics course (though we use statistics). It is not a computer science course (though we write code). It is a \emph{practical data analysis course} for people whose primary expertise is health, not programming.

\paragraph{Prerequisites.} None beyond basic computer literacy. Familiarity with descriptive statistics (mean, median, standard deviation) is helpful but will be reviewed.

\paragraph{Software.} All code runs in Jupyter notebooks, either locally (Anaconda) or in the cloud (Google Colab --- free, no installation required).

\vfill
\noindent Ya\'{e} Ulrich Gaba\\
Cotonou, 2026

% ============================================================
\mainmatter

\chapter{Python essentials for health professionals}
\label{ch:python-essentials}

\begin{quote}
\emph{``The goal is not to learn Python. The goal is to answer health questions with data. Python is just the tool.''}
\end{quote}

% ============================================================
\section{Why Python for health data?}

Three reasons health professionals should learn Python rather than relying solely on Excel or SPSS:

\begin{enumerate}
  \item \textbf{Reproducibility.} A Python script documents every step of your analysis. When a reviewer asks ``how did you handle missing lab values?'', you show the code --- not a description of menu clicks.
  \item \textbf{Scale.} Excel struggles with 100\,000 patient records. Python handles millions without slowing down.
  \item \textbf{Ecosystem.} Libraries like \texttt{lifelines} (survival analysis), \texttt{geopandas} (disease mapping), and \texttt{scikit-learn} (predictive models) give you tools that SPSS does not offer.
\end{enumerate}

\begin{clinicalnote}
You will not need to memorize syntax. Health professionals use Python like clinicians use a stethoscope --- you learn the parts you need, and you look up the rest. Every code block in this course can be copied into a Jupyter notebook and run immediately.
\end{clinicalnote}

% ============================================================
\section{Setting up your environment}

\subsection{Google Colab (recommended)}

No installation. Go to \url{https://colab.research.google.com}, sign in with a Google account, and create a new notebook. Every library we use is pre-installed.

\subsection{Local installation with Anaconda}

If you prefer working offline:
\begin{enumerate}
  \item Download Anaconda: \url{https://www.anaconda.com/download}
  \item Install with default settings.
  \item Open \textbf{Jupyter Notebook} from Anaconda Navigator.
\end{enumerate}

% ============================================================
\section{Your first notebook}

A Jupyter notebook is a document that mixes text, code, and output. Each \emph{cell} contains either:
\begin{itemize}
  \item \textbf{Code} --- Python instructions that produce a result.
  \item \textbf{Markdown} --- formatted text (explanations, headings, notes).
\end{itemize}

Press \texttt{Shift + Enter} to run a cell.

\begin{minted}{python}
# Your first Python cell
print("Hello, health data!")
\end{minted}

% ============================================================
\section{Variables and types}

A \emph{variable} stores a value. In health data, you work with four types constantly:

\begin{minted}{python}
patient_age = 45              # int (integer)
temperature = 37.2            # float (decimal number)
diagnosis = "Type 2 Diabetes" # str (string / text)
is_hospitalized = True        # bool (True or False)
\end{minted}

\begin{pythontip}
Variable names should be descriptive. Use \texttt{patient\_age}, not \texttt{x}. Your future self (and your collaborators) will thank you.
\end{pythontip}

% ============================================================
\section{Collections: lists and dictionaries}

\subsection{Lists}

A list holds multiple values in order:
\begin{minted}{python}
blood_pressures = [120, 135, 128, 142, 118]
print(f"Number of readings: {len(blood_pressures)}")
print(f"First reading: {blood_pressures[0]}")  # indexing starts at 0
print(f"Last reading: {blood_pressures[-1]}")
\end{minted}

\subsection{Dictionaries}

A dictionary maps keys to values --- like a patient chart:
\begin{minted}{python}
patient = {
    "id": "PAT-0042",
    "age": 58,
    "sex": "F",
    "diagnosis": "Hypertension",
    "systolic_bp": 148,
    "medications": ["Amlodipine", "Hydrochlorothiazide"]
}
print(patient["diagnosis"])  # Hypertension
\end{minted}

% ============================================================
\section{Control flow}

\subsection{Conditionals}

\begin{minted}{python}
systolic = 148

if systolic >= 140:
    category = "Stage 2 Hypertension"
elif systolic >= 130:
    category = "Stage 1 Hypertension"
elif systolic >= 120:
    category = "Elevated"
else:
    category = "Normal"

print(f"BP category: {category}")
\end{minted}

\subsection{Loops}

\begin{minted}{python}
patients = [
    {"id": "P001", "age": 34, "hba1c": 5.4},
    {"id": "P002", "age": 67, "hba1c": 7.8},
    {"id": "P003", "age": 52, "hba1c": 6.1},
]

for p in patients:
    if p["hba1c"] >= 6.5:
        print(f"{p['id']}: Diabetic (HbA1c = {p['hba1c']})")
    elif p["hba1c"] >= 5.7:
        print(f"{p['id']}: Pre-diabetic (HbA1c = {p['hba1c']})")
    else:
        print(f"{p['id']}: Normal (HbA1c = {p['hba1c']})")
\end{minted}

% ============================================================
\section{Functions}

A function encapsulates a reusable computation:

\begin{minted}{python}
def bmi(weight_kg, height_m):
    """Calculate Body Mass Index."""
    return weight_kg / (height_m ** 2)

def bmi_category(bmi_value):
    """Classify BMI according to WHO categories."""
    if bmi_value < 18.5:
        return "Underweight"
    elif bmi_value < 25.0:
        return "Normal weight"
    elif bmi_value < 30.0:
        return "Overweight"
    else:
        return "Obese"

value = bmi(82, 1.75)
print(f"BMI: {value:.1f} ({bmi_category(value)})")
\end{minted}

\begin{exercise}
\begin{enumerate}
  \item Write a function \texttt{classify\_bp(systolic, diastolic)} that returns a blood pressure category according to the AHA/ACC guidelines.
  \item Create a list of 5 patient dictionaries with \texttt{weight\_kg} and \texttt{height\_m} fields. Loop through them, compute BMI, and print the category for each.
  \item Write a function \texttt{egfr(creatinine, age, sex)} that estimates glomerular filtration rate using the CKD-EPI equation.
\end{enumerate}
\end{exercise}

% ============================================================
\section{Importing libraries}

Python's power comes from libraries. Here are the ones we will use throughout this course:

\begin{minted}{python}
import pandas as pd          # data manipulation
import numpy as np           # numerical computing
import matplotlib.pyplot as plt  # plotting
import seaborn as sns        # statistical visualization
from scipy import stats      # statistical tests
\end{minted}

\begin{warningbox}
If you get \texttt{ModuleNotFoundError}, install the missing library:
\begin{minted}{bash}
pip install pandas matplotlib seaborn scipy
\end{minted}
In Google Colab, prefix with \texttt{!}: \texttt{!pip install lifelines}
\end{warningbox}

% ============================================================
\section{Chapter summary}

\begin{itemize}
  \item Python is free, reproducible, and has health-specific libraries.
  \item Jupyter notebooks mix code, text, and output in one document.
  \item Variables store values; lists and dictionaries organize them.
  \item Functions make your analysis reusable and readable.
  \item \texttt{pandas}, \texttt{matplotlib}, \texttt{seaborn}, and \texttt{scipy} are the core libraries for health data analysis.
\end{itemize}

\chapter{Health data with Pandas}
\label{ch:pandas}

\begin{quote}
\emph{``80\% of data analysis is getting the data into a shape where you can ask it questions.''}
\end{quote}

% ============================================================
\section{What is a DataFrame?}

A DataFrame is a table --- rows are observations (patients, countries, time points), columns are variables (age, diagnosis, lab values). If you have used Excel, you already understand the concept. Pandas gives you Excel's power with Python's reproducibility.

\begin{minted}{python}
import pandas as pd

# Create a small clinical DataFrame
data = {
    "patient_id": ["P001", "P002", "P003", "P004", "P005"],
    "age": [45, 67, 34, 52, 71],
    "sex": ["M", "F", "F", "M", "F"],
    "systolic_bp": [128, 158, 112, 145, 162],
    "hba1c": [5.4, 7.8, 5.1, 6.3, 8.2],
    "diagnosis": ["None", "T2DM", "None", "Pre-DM", "T2DM"]
}
df = pd.DataFrame(data)
df
\end{minted}

% ============================================================
\section{Loading real health data}

\subsection{CSV files}

Most health datasets come as CSV (comma-separated values) files:

\begin{minted}{python}
# WHO life expectancy data (Gapminder)
url = "https://raw.githubusercontent.com/datasets/gapminder/main/data/gapminder.csv"
gapminder = pd.read_csv(url)
print(f"Shape: {gapminder.shape}")  # (rows, columns)
gapminder.head()
\end{minted}

\subsection{Excel files}

\begin{minted}{python}
df = pd.read_excel("patient_records.xlsx", sheet_name="Lab Results")
\end{minted}

\subsection{From the WHO Global Health Observatory}

\begin{minted}{python}
# Maternal mortality ratio by country
url = "https://apps.who.int/gho/athena/api/GHO/MDG_0000000026?format=csv"
mmr = pd.read_csv(url)
mmr.head()
\end{minted}

\begin{clinicalnote}
The WHO GHO API provides over 1\,000 health indicators for 194 member states. You can browse indicators at \url{https://www.who.int/data/gho/data/indicators}. Each indicator has a code (e.g., \texttt{MDG\_0000000026} for maternal mortality ratio) that you plug into the API URL.
\end{clinicalnote}

% ============================================================
\section{Exploring your data}

The first thing you do with any health dataset --- before any analysis --- is \emph{look at it}.

\begin{minted}{python}
# Basic exploration
print(gapminder.shape)          # dimensions
print(gapminder.columns.tolist())  # column names
print(gapminder.dtypes)         # data types
gapminder.describe()            # summary statistics
gapminder.info()                # non-null counts and memory
\end{minted}

\subsection{Selecting columns}

\begin{minted}{python}
# Single column (returns a Series)
ages = gapminder["lifeExp"]

# Multiple columns (returns a DataFrame)
subset = gapminder[["country", "year", "lifeExp"]]
\end{minted}

\subsection{Filtering rows}

\begin{minted}{python}
# Patients with HbA1c >= 6.5 (diabetic threshold)
diabetic = df[df["hba1c"] >= 6.5]

# African countries in Gapminder
africa = gapminder[gapminder["continent"] == "Africa"]

# Multiple conditions: female patients over 60 with high BP
high_risk = df[(df["sex"] == "F") & (df["age"] > 60) & (df["systolic_bp"] >= 140)]
\end{minted}

\begin{warningbox}
Use \texttt{\&} (and), \texttt{|} (or), \texttt{\textasciitilde} (not) with parentheses around each condition. Python's \texttt{and}/\texttt{or} keywords do not work with DataFrames.
\end{warningbox}

% ============================================================
\section{Sorting and ranking}

\begin{minted}{python}
# Countries with highest life expectancy in 2007
recent = gapminder[gapminder["year"] == 2007]
top10 = recent.sort_values("lifeExp", ascending=False).head(10)
print(top10[["country", "lifeExp"]])
\end{minted}

% ============================================================
\section{Grouping and aggregation}

Grouping is how you answer questions like ``what is the average life expectancy per continent?'' or ``how many patients in each diagnosis category?''

\begin{minted}{python}
# Average life expectancy by continent (2007)
recent.groupby("continent")["lifeExp"].mean().sort_values(ascending=False)
\end{minted}

\begin{minted}{python}
# Multiple aggregations
recent.groupby("continent")["lifeExp"].agg(["mean", "median", "std", "count"])
\end{minted}

\begin{minted}{python}
# Patient counts by diagnosis
df.groupby("diagnosis")["patient_id"].count()
\end{minted}

% ============================================================
\section{Creating new columns}

\begin{minted}{python}
# BMI from weight and height
df["bmi"] = df["weight_kg"] / (df["height_m"] ** 2)

# Age group
df["age_group"] = pd.cut(df["age"],
                         bins=[0, 18, 40, 60, 100],
                         labels=["<18", "18-39", "40-59", "60+"])

# Binary flag
df["hypertensive"] = (df["systolic_bp"] >= 140).astype(int)
\end{minted}

% ============================================================
\section{Saving your results}

\begin{minted}{python}
# Save to CSV
df.to_csv("cleaned_patients.csv", index=False)

# Save to Excel
df.to_excel("cleaned_patients.xlsx", index=False, sheet_name="Patients")
\end{minted}

% ============================================================
\section{Practical: WHO life expectancy analysis}

\begin{exercise}
Using the Gapminder dataset:
\begin{enumerate}
  \item Load the data and display basic statistics.
  \item Filter to African countries only. How many unique countries?
  \item Compute mean life expectancy by decade (1950s, 1960s, ..., 2000s) for Africa vs.\ Europe.
  \item Find the 5 African countries with the lowest life expectancy in 2007.
  \item Create a column \texttt{life\_exp\_category}: ``Low'' ($<55$), ``Medium'' ($55$--$70$), ``High'' ($>70$).
  \item Group by continent and \texttt{life\_exp\_category}. How many countries in each?
  \item Save the Africa-only data to a CSV file.
\end{enumerate}
\end{exercise}

% ============================================================
\section{Chapter summary}

\begin{itemize}
  \item A DataFrame is a table. Rows = observations, columns = variables.
  \item \texttt{pd.read\_csv()} loads data; \texttt{.head()}, \texttt{.describe()}, \texttt{.info()} explore it.
  \item Filter with boolean conditions: \texttt{df[df["col"] > value]}.
  \item \texttt{.groupby()} + \texttt{.agg()} answers ``per group'' questions.
  \item \texttt{pd.cut()} creates categorical variables from continuous ones.
  \item Always save cleaned data for reproducibility.
\end{itemize}

\chapter{Data cleaning in health contexts}
\label{ch:data-cleaning}

\begin{quote}
\emph{``Garbage in, garbage out. In health data, garbage in can mean wrong treatments, wrong policies, and wrong conclusions about human lives.''}
\end{quote}

% ============================================================
\section{Why data cleaning matters in health}

Every health dataset has problems. Lab values get entered with wrong units. Patients appear twice under different IDs. Dates are recorded as \texttt{03/04/2023} --- is that March 4 or April 3? A blood glucose of 5000~mg/dL is clearly an error, but a glucose of 500~mg/dL might be a real diabetic emergency.

Data cleaning is not optional pre-processing. It is a clinical reasoning task. You must decide:
\begin{itemize}
  \item Is this value \emph{wrong} or just \emph{unusual}?
  \item Should this record be \emph{corrected}, \emph{flagged}, or \emph{removed}?
  \item Will my cleaning decisions \emph{bias} the analysis?
\end{itemize}

\begin{clinicalnote}
In a 2019 study of electronic health records from a large US hospital system, researchers found that 18\% of lab values had at least one quality issue --- duplicate entries, physiologically impossible values, or inconsistent units. Cleaning decisions changed the estimated prevalence of chronic kidney disease by over 3 percentage points.
\end{clinicalnote}

% ============================================================
\section{Missing data: types and mechanisms}

\subsection{The three mechanisms}

Not all missing data is the same. The mechanism determines what you can do about it.

\begin{enumerate}
  \item \textbf{MCAR --- Missing Completely At Random.} The missingness has nothing to do with the data. Example: a lab tube is dropped and the sample is lost. The probability of losing the sample is unrelated to the patient's health.
  \item \textbf{MAR --- Missing At Random.} The missingness depends on \emph{observed} variables but not the missing value itself. Example: younger patients are less likely to have cholesterol measured (because clinicians order it less often for young patients). Once you account for age, the missingness is random.
  \item \textbf{MNAR --- Missing Not At Random.} The missingness depends on the \emph{unobserved} value itself. Example: very sick patients are too ill to attend follow-up appointments, so their outcome data is missing \emph{because} they are sick. This is the hardest case.
\end{enumerate}

\begin{warningbox}
MNAR is common in health data and dangerous. If you drop all patients with missing follow-up data, you are systematically excluding the sickest patients --- your analysis will overestimate treatment success. There is no statistical fix for MNAR; you need domain knowledge to assess the likely direction and magnitude of bias.
\end{warningbox}

\subsection{Clinical examples of each mechanism}

\begin{center}
\begin{tabular}{lp{5cm}p{5cm}}
\toprule
\textbf{Type} & \textbf{Example} & \textbf{Consequence of dropping} \\
\midrule
MCAR & Lab sample hemolyzed due to transport vibration & No bias, but reduced sample size \\
MAR & HbA1c not ordered for non-diabetic patients & Bias if you study HbA1c without accounting for indication \\
MNAR & Severely depressed patients skip follow-up questionnaires & Underestimate depression severity in remaining data \\
\bottomrule
\end{tabular}
\end{center}

% ============================================================
\section{Detecting missing values with pandas}

\subsection{Basic detection}

\begin{minted}{python}
import pandas as pd
import numpy as np

# Create a clinical dataset with missing values
data = {
    "patient_id": ["P001", "P002", "P003", "P004", "P005", "P006"],
    "age": [45, 67, 34, np.nan, 71, 52],
    "sex": ["M", "F", "F", "M", np.nan, "F"],
    "systolic_bp": [128, 158, np.nan, 145, 162, np.nan],
    "hba1c": [5.4, 7.8, 5.1, np.nan, 8.2, 6.3],
    "smoking": ["Never", np.nan, "Current", "Former", np.nan, "Never"]
}
df = pd.DataFrame(data)

# Count missing values per column
print(df.isnull().sum())

# Percentage missing per column
print((df.isnull().sum() / len(df) * 100).round(1))

# Rows with any missing value
print(f"Rows with missing data: {df.isnull().any(axis=1).sum()} / {len(df)}")
\end{minted}

\begin{pythontip}
\texttt{.isnull()} and \texttt{.isna()} are identical in pandas. Use whichever reads better to you. Both detect \texttt{NaN}, \texttt{None}, and \texttt{NaT} (missing datetime).
\end{pythontip}

\subsection{Visualizing missingness with msno}

The \texttt{missingno} library provides instant visual summaries of missing data patterns:

\begin{minted}{python}
import missingno as msno
import matplotlib.pyplot as plt

# Matrix plot: white bars = missing
msno.matrix(df)
plt.title("Missing data pattern")
plt.tight_layout()
plt.show()

# Bar chart: count of non-null values per column
msno.bar(df)
plt.tight_layout()
plt.show()

# Heatmap: correlations between missingness of different columns
# If two columns tend to be missing together, they will show high correlation
msno.heatmap(df)
plt.tight_layout()
plt.show()
\end{minted}

\begin{warningbox}
If \texttt{missingno} is not installed, run \texttt{!pip install missingno} in your notebook. In Google Colab it is not pre-installed.
\end{warningbox}

\subsection{Patterns of co-missingness}

When two variables are often missing together, it suggests a common cause. For example, if both \texttt{systolic\_bp} and \texttt{diastolic\_bp} are missing for the same patients, it likely means the BP measurement was not taken at all --- not that each value was independently lost.

\begin{minted}{python}
# Check which columns are missing together
missing_pairs = df.isnull().corr()
print(missing_pairs)
\end{minted}

% ============================================================
\section{Imputation strategies}

\subsection{When to impute vs.\ when to drop}

\begin{itemize}
  \item \textbf{Drop rows} if missingness is MCAR and the missing fraction is small ($<5\%$).
  \item \textbf{Drop columns} if $>50\%$ of values are missing and the variable is not critical.
  \item \textbf{Impute} when you need to retain sample size and you have a reasonable strategy.
  \item \textbf{Flag + sensitivity analysis} for MNAR data.
\end{itemize}

\begin{minted}{python}
# Drop rows with any missing value
df_complete = df.dropna()
print(f"Rows remaining: {len(df_complete)} / {len(df)}")

# Drop rows only if specific columns are missing
df_partial = df.dropna(subset=["age", "systolic_bp"])
\end{minted}

\subsection{Mean/median imputation for continuous variables}

\begin{minted}{python}
# Median imputation for lab values (robust to outliers)
df["systolic_bp"] = df["systolic_bp"].fillna(df["systolic_bp"].median())
df["hba1c"] = df["hba1c"].fillna(df["hba1c"].median())

# Mean imputation for age (approximately symmetric distribution)
df["age"] = df["age"].fillna(df["age"].mean())
\end{minted}

\begin{clinicalnote}
Use \textbf{median} rather than mean for lab values like creatinine, glucose, and triglycerides. These distributions are typically right-skewed --- a few very high values pull the mean upward, making it a poor representative of the ``typical'' patient. The median is resistant to these outliers.
\end{clinicalnote}

\subsection{Mode imputation for categorical variables}

\begin{minted}{python}
# Mode imputation for sex and smoking status
df["sex"] = df["sex"].fillna(df["sex"].mode()[0])
df["smoking"] = df["smoking"].fillna(df["smoking"].mode()[0])
\end{minted}

\subsection{Group-specific imputation}

Sometimes the right imputation value depends on a group. Imputing systolic BP with the overall median ignores the fact that older patients have higher BP:

\begin{minted}{python}
# Impute systolic_bp with the median for the patient's age group
df["age_group"] = pd.cut(df["age"], bins=[0, 40, 60, 100],
                         labels=["<40", "40-59", "60+"])
df["systolic_bp"] = df.groupby("age_group")["systolic_bp"].transform(
    lambda x: x.fillna(x.median())
)
\end{minted}

\subsection{KNN imputation}

K-nearest neighbors imputation finds the $k$ most similar patients (based on other variables) and uses their values to fill in the missing one:

\begin{minted}{python}
from sklearn.impute import KNNImputer

# Select numeric columns for KNN imputation
numeric_cols = ["age", "systolic_bp", "hba1c"]
imputer = KNNImputer(n_neighbors=5)
df[numeric_cols] = imputer.fit_transform(df[numeric_cols])
\end{minted}

\begin{pythontip}
KNN imputation respects relationships between variables. If a patient with high BMI, high age, and high HbA1c is missing systolic BP, KNN will impute from similar patients --- who likely also have high BP. This is more realistic than using the overall median.
\end{pythontip}

\subsection{Creating a missingness indicator}

For important variables, create a flag so you can later test whether missingness itself is associated with outcomes:

\begin{minted}{python}
# Before imputing, create a flag
df["bp_was_missing"] = df["systolic_bp"].isnull().astype(int)

# Then impute
df["systolic_bp"] = df["systolic_bp"].fillna(df["systolic_bp"].median())
\end{minted}

% ============================================================
\section{Handling duplicates}

\subsection{Detecting duplicates}

\begin{minted}{python}
# Load a dataset with deliberate duplicates
patients = pd.DataFrame({
    "patient_id": ["P001", "P002", "P003", "P002", "P004", "P003"],
    "visit_date": ["2023-01-15", "2023-01-16", "2023-01-17",
                   "2023-01-16", "2023-01-18", "2023-02-20"],
    "systolic_bp": [128, 158, 112, 158, 145, 118],
    "diagnosis": ["HTN", "T2DM", "None", "T2DM", "HTN", "None"]
})

# Exact duplicates (all columns identical)
print(f"Exact duplicates: {patients.duplicated().sum()}")
print(patients[patients.duplicated(keep=False)])  # show all copies

# Duplicates based on patient_id only
print(f"\nDuplicate patient IDs: {patients.duplicated(subset='patient_id').sum()}")
\end{minted}

\subsection{Deciding what counts as a duplicate}

In health data, ``duplicate'' is not always straightforward:

\begin{center}
\begin{tabular}{lp{8cm}}
\toprule
\textbf{Scenario} & \textbf{Appropriate action} \\
\midrule
Same patient, same visit, identical values & True duplicate --- remove one copy \\
Same patient, same visit, different values & Data entry error --- investigate which is correct \\
Same patient, different visit dates & Repeated measures --- keep all (this is longitudinal data) \\
Different patient IDs, same demographics & Possible duplicate registration --- flag for manual review \\
\bottomrule
\end{tabular}
\end{center}

\begin{minted}{python}
# Remove exact duplicates
patients_clean = patients.drop_duplicates()
print(f"After removing exact duplicates: {len(patients_clean)} rows")

# Keep only the first visit per patient
first_visits = patients.drop_duplicates(subset="patient_id", keep="first")
print(f"First visits only: {len(first_visits)} rows")

# Keep only the last visit per patient
last_visits = patients.drop_duplicates(subset="patient_id", keep="last")
\end{minted}

\begin{clinicalnote}
In clinical registries, duplicate patient records are a major safety issue. A patient registered under two IDs may receive conflicting prescriptions. Deduplication in production systems uses probabilistic matching on name, date of birth, and address --- far beyond what simple \texttt{drop\_duplicates()} can do.
\end{clinicalnote}

% ============================================================
\section{Inconsistent data}

\subsection{ICD code formats}

The International Classification of Diseases uses codes like \texttt{E11.9} (Type 2 diabetes without complications). In messy data, you will find the same condition recorded as \texttt{E11.9}, \texttt{E119}, \texttt{e11.9}, or even \texttt{250.00} (the old ICD-9 code).

\begin{minted}{python}
diagnoses = pd.Series(["E11.9", "e11.9", "E119", "E11.9 ", " e11.9"])

# Standardize: uppercase, strip whitespace, ensure dot format
cleaned = (diagnoses
           .str.strip()
           .str.upper()
           .str.replace(r"^([A-Z]\d{2})(\d+)$", r"\1.\2", regex=True))
print(cleaned)
\end{minted}

\subsection{Date formats}

\begin{minted}{python}
dates = pd.Series(["2023-01-15", "01/15/2023", "15-Jan-2023",
                   "Jan 15, 2023", "20230115"])

# pd.to_datetime handles multiple formats automatically
parsed = pd.to_datetime(dates, format="mixed", dayfirst=False)
print(parsed)

# Standardize to ISO format
print(parsed.dt.strftime("%Y-%m-%d"))
\end{minted}

\begin{warningbox}
The date \texttt{03/04/2023} is March 4 in the US but April 3 in Europe. Always check the source. If your dataset mixes conventions, you will get silent errors --- the date will parse without warning, but it will be \emph{wrong}. Use \texttt{dayfirst=True} for European-format dates.
\end{warningbox}

\subsection{Unit conversions}

Different labs report results in different units. The most common confusion:

\begin{center}
\begin{tabular}{llll}
\toprule
\textbf{Analyte} & \textbf{US unit} & \textbf{SI unit} & \textbf{Conversion} \\
\midrule
Glucose & mg/dL & mmol/L & divide by 18.018 \\
Cholesterol & mg/dL & mmol/L & divide by 38.67 \\
Creatinine & mg/dL & \textmu mol/L & multiply by 88.42 \\
Hemoglobin & g/dL & g/L & multiply by 10 \\
\bottomrule
\end{tabular}
\end{center}

\begin{minted}{python}
# Standardize glucose to mmol/L
# Assume values > 30 are in mg/dL (glucose in mmol/L is rarely > 30)
def standardize_glucose(value, unit=None):
    """Convert glucose to mmol/L. Infer unit if not provided."""
    if unit == "mg/dL" or (unit is None and value > 30):
        return value / 18.018
    return value  # already in mmol/L

df["glucose_mmol"] = df["glucose"].apply(
    lambda x: standardize_glucose(x)
)
\end{minted}

\begin{pythontip}
When a dataset has no unit column, use the value range to infer the unit. Fasting glucose in mg/dL is typically 70--300; in mmol/L it is 3.9--16.7. If you see a value of 126, it is almost certainly mg/dL. If you see 7.0, it is almost certainly mmol/L.
\end{pythontip}

\subsection{Standardizing text fields}

\begin{minted}{python}
# Sex/gender field with inconsistent entries
sex_col = pd.Series(["Male", "male", "M", "m", "MALE", "Female",
                     "female", "F", "f", "FEMALE"])

# Map to standard values
sex_mapping = {"male": "M", "m": "M", "female": "F", "f": "F"}
sex_clean = sex_col.str.strip().str.lower().map(sex_mapping)
print(sex_clean)
\end{minted}

% ============================================================
\section{Data validation}

After cleaning, validate that your data makes clinical sense.

\subsection{Range checks}

\begin{minted}{python}
# Define clinically plausible ranges
ranges = {
    "age": (0, 120),
    "systolic_bp": (60, 300),
    "diastolic_bp": (30, 200),
    "heart_rate": (20, 300),
    "hba1c": (2.0, 20.0),
    "glucose_mgdl": (10, 1500),
    "bmi": (8, 80),
    "temperature_c": (25, 45),
}

def validate_range(df, column, low, high):
    """Flag values outside the plausible range."""
    out_of_range = df[(df[column] < low) | (df[column] > high)]
    if len(out_of_range) > 0:
        print(f"WARNING: {len(out_of_range)} values in '{column}' "
              f"outside [{low}, {high}]")
        print(out_of_range[["patient_id", column]])
    return out_of_range

# Apply range checks
for col, (low, high) in ranges.items():
    if col in df.columns:
        validate_range(df, col, low, high)
\end{minted}

\subsection{Cross-field validation}

Some errors only become visible when you compare columns:

\begin{minted}{python}
# Diastolic should be lower than systolic
bp_errors = df[df["diastolic_bp"] >= df["systolic_bp"]]
if len(bp_errors) > 0:
    print(f"WARNING: {len(bp_errors)} rows where diastolic >= systolic")

# Pregnancy flag should only be 1 for female patients
preg_errors = df[(df["pregnant"] == 1) & (df["sex"] == "M")]
if len(preg_errors) > 0:
    print(f"WARNING: {len(preg_errors)} male patients flagged as pregnant")

# Death date should be after admission date
date_errors = df[df["death_date"] < df["admission_date"]]
if len(date_errors) > 0:
    print(f"WARNING: {len(date_errors)} deaths recorded before admission")
\end{minted}

\begin{clinicalnote}
Not every ``impossible'' value is an error. A heart rate of 250~bpm is physiologically extreme but can occur in supraventricular tachycardia. A glucose of 800~mg/dL is rare but real in diabetic ketoacidosis. Domain knowledge tells you whether to remove, flag, or keep.
\end{clinicalnote}

\subsection{Automated validation report}

\begin{minted}{python}
def data_quality_report(df):
    """Generate a summary of data quality issues."""
    report = []
    for col in df.columns:
        n_missing = df[col].isnull().sum()
        pct_missing = n_missing / len(df) * 100
        n_unique = df[col].nunique()

        row = {
            "column": col,
            "dtype": str(df[col].dtype),
            "n_missing": n_missing,
            "pct_missing": round(pct_missing, 1),
            "n_unique": n_unique,
        }

        if df[col].dtype in ["float64", "int64"]:
            row["min"] = df[col].min()
            row["max"] = df[col].max()
            row["mean"] = round(df[col].mean(), 2)

        report.append(row)

    return pd.DataFrame(report)

quality = data_quality_report(df)
print(quality.to_string(index=False))
\end{minted}

% ============================================================
\section{A complete cleaning pipeline}

Here is a realistic pipeline applied to NHANES-style data:

\begin{minted}{python}
import pandas as pd
import numpy as np
from sklearn.impute import KNNImputer

# -----------------------------------------------------------
# Step 1: Load data
# -----------------------------------------------------------
url = ("https://raw.githubusercontent.com/datasets/"
       "gapminder/main/data/gapminder.csv")
df = pd.read_csv(url)
print(f"Raw data: {df.shape}")

# -----------------------------------------------------------
# Step 2: Standardize column names
# -----------------------------------------------------------
df.columns = df.columns.str.strip().str.lower().str.replace(" ", "_")

# -----------------------------------------------------------
# Step 3: Remove exact duplicates
# -----------------------------------------------------------
n_before = len(df)
df = df.drop_duplicates()
print(f"Duplicates removed: {n_before - len(df)}")

# -----------------------------------------------------------
# Step 4: Assess missing data
# -----------------------------------------------------------
missing_pct = (df.isnull().sum() / len(df) * 100).round(1)
print("Missing percentages:")
print(missing_pct[missing_pct > 0])

# -----------------------------------------------------------
# Step 5: Validate ranges
# -----------------------------------------------------------
# Life expectancy should be between 20 and 90
outliers = df[(df["lifeexp"] < 20) | (df["lifeexp"] > 90)]
print(f"Life expectancy outliers: {len(outliers)}")

# -----------------------------------------------------------
# Step 6: Save cleaned data
# -----------------------------------------------------------
df.to_csv("gapminder_cleaned.csv", index=False)
print(f"Cleaned data saved: {df.shape}")
\end{minted}

% ============================================================
\section{Practical: cleaning a messy clinical dataset}

\begin{exercise}
Download or create a messy clinical dataset with the following deliberate errors, then clean it step by step:

\begin{minted}{python}
import pandas as pd
import numpy as np

# Messy clinical data with deliberate errors
messy = pd.DataFrame({
    "patient_id": ["P001", "P002", "P003", "P004", "P002",
                   "P005", "P006", "P007", "P008", "P009"],
    "age": [45, 67, -3, 52, 67, 200, 38, np.nan, 71, 44],
    "sex": ["M", "F", "Female", "m", "F",
            "Male", np.nan, "F", "X", "M"],
    "systolic_bp": [128, 158, 112, 450, 158,
                    135, 122, np.nan, 162, 118],
    "diastolic_bp": [82, 160, 74, 92, 160,
                     88, 78, np.nan, 98, 76],
    "glucose_mgdl": [95, 7.2, 110, 180, 95,
                     88, 102, 6.8, 220, 5.5],
    "hba1c": [5.4, 7.8, 5.1, 6.3, 7.8,
              np.nan, 5.8, 6.1, 55, 5.2],
    "visit_date": ["2023-01-15", "01/16/2023", "2023-01-17",
                   "2023/01/18", "2023-01-16",
                   "15-Jan-2023", "2023-01-20", "2023-01-21",
                   "2023-01-22", "Jan 23, 2023"],
    "icd_code": ["I10", "E119", "e11.9", "I10 ", " i10",
                 "E11.9", "I10", "e119", "I10", "E11.9"]
})
\end{minted}

Tasks:
\begin{enumerate}
  \item Generate a data quality report: count missing values, check data types, compute summary statistics.
  \item Identify and remove exact duplicate rows (P002 appears twice with identical data).
  \item Fix the \texttt{sex} column: standardize to ``M'' and ``F''. Decide what to do with ``X'' and \texttt{NaN}.
  \item Validate \texttt{age}: flag or correct the impossible values ($-3$ and $200$). Set them to \texttt{NaN}.
  \item Validate \texttt{systolic\_bp}: a value of 450 is impossible. Set to \texttt{NaN}.
  \item Check that diastolic $<$ systolic for all patients. Flag violations.
  \item Fix \texttt{glucose\_mgdl}: some values (7.2, 6.8, 5.5) are clearly in mmol/L. Convert them to mg/dL by multiplying by 18.018.
  \item Fix \texttt{hba1c}: the value 55 is likely in mmol/mol (IFCC) rather than \% (NGSP). Convert using: $\text{NGSP} = 0.0915 \times \text{IFCC} + 2.15$.
  \item Standardize \texttt{visit\_date} to ISO format (YYYY-MM-DD).
  \item Standardize \texttt{icd\_code}: uppercase, stripped, with dot separator.
  \item Impute remaining missing values: median for continuous, mode for categorical.
  \item Generate a final data quality report and compare with the initial one.
  \item Save the cleaned dataset to \texttt{cleaned\_clinical\_data.csv}.
\end{enumerate}
\end{exercise}

% ============================================================
\section{Chapter summary}

\begin{itemize}
  \item Missing data in health has three mechanisms: MCAR, MAR, and MNAR. The mechanism determines the appropriate strategy.
  \item Use \texttt{.isnull().sum()} for quick counts and \texttt{missingno} for visual patterns.
  \item Imputation options: mean/median for continuous data, mode for categorical, KNN for multivariate imputation.
  \item Duplicates in health data may be true duplicates, data entry errors, or repeated visits --- each requires a different response.
  \item Inconsistent coding (ICD formats, date formats, units) is ubiquitous and must be standardized.
  \item Validate data against clinically plausible ranges \emph{after} cleaning.
  \item Always create a data quality report before and after cleaning so you can document every decision.
\end{itemize}

\chapter{Descriptive statistics and epidemiological measures}
\label{ch:descriptive-epi}

\begin{quote}
\emph{``Before you fit a model, describe the data. Before you describe the data, look at the data. Epidemiology begins with counting.''}
\end{quote}

% ============================================================
\section{Central tendency: what is ``typical''?}

\subsection{Mean, median, and when to use which}

\begin{minted}{python}
import pandas as pd
import numpy as np

# Fasting blood glucose for 10 patients (mg/dL)
glucose = pd.Series([92, 98, 104, 95, 88, 110, 97, 340, 101, 93])

print(f"Mean:   {glucose.mean():.1f} mg/dL")
print(f"Median: {glucose.median():.1f} mg/dL")
\end{minted}

The mean is 121.8~mg/dL. The median is 97.5~mg/dL. One patient with a glucose of 340 (diabetic ketoacidosis) pulls the mean up by nearly 25~mg/dL. The median is unaffected.

\begin{clinicalnote}
\textbf{Rule of thumb for health data:} use the \textbf{median} when the distribution is skewed (most lab values, length of stay, health expenditures). Use the \textbf{mean} when the distribution is approximately symmetric (height, diastolic blood pressure in healthy adults).
\end{clinicalnote}

\subsection{Mode}

The mode is the most frequent value. It is the only measure of central tendency for categorical data:

\begin{minted}{python}
diagnoses = pd.Series(["HTN", "T2DM", "HTN", "None", "HTN",
                       "T2DM", "None", "HTN", "COPD", "HTN"])
print(f"Most common diagnosis: {diagnoses.mode()[0]}")
print(f"Frequency: {diagnoses.value_counts().iloc[0]}/{len(diagnoses)}")
\end{minted}

% ============================================================
\section{Spread: how variable is the data?}

\subsection{Standard deviation and variance}

\begin{minted}{python}
# Blood pressure readings from two clinics
clinic_a = pd.Series([120, 122, 118, 125, 121, 119, 123, 120])
clinic_b = pd.Series([98, 145, 110, 155, 102, 160, 115, 135])

print(f"Clinic A: mean={clinic_a.mean():.1f}, SD={clinic_a.std():.1f}")
print(f"Clinic B: mean={clinic_b.mean():.1f}, SD={clinic_b.std():.1f}")
\end{minted}

Both clinics have similar means ($\approx$121), but Clinic~B has much higher variability. A physician looking only at the mean would miss that Clinic~B has patients with dangerously high \emph{and} unusually low blood pressure.

\subsection{IQR and percentiles}

The interquartile range (IQR) is the range between the 25th and 75th percentiles. Like the median, it is robust to outliers.

\begin{minted}{python}
# Length of hospital stay (days) --- typically right-skewed
los = pd.Series([2, 3, 3, 4, 4, 5, 5, 6, 7, 8, 12, 28, 45])

q25 = los.quantile(0.25)
q75 = los.quantile(0.75)
iqr = q75 - q25

print(f"Median LOS: {los.median():.0f} days")
print(f"IQR: {q25:.0f}--{q75:.0f} days (range: {iqr:.0f})")
print(f"Mean LOS: {los.mean():.1f} days")  # inflated by 28 and 45
\end{minted}

\begin{pythontip}
Pandas \texttt{.describe()} gives you count, mean, std, min, 25\%, 50\%, 75\%, and max in one call. It is the fastest way to get a statistical snapshot of every numeric column.
\end{pythontip}

\subsection{Coefficient of variation}

The coefficient of variation (CV) expresses the standard deviation as a percentage of the mean. It is useful for comparing variability between measurements with different units or scales:

\begin{minted}{python}
# Compare variability of glucose (mg/dL) vs hemoglobin (g/dL)
glucose = pd.Series([95, 102, 88, 110, 97, 105, 92, 98])
hemoglobin = pd.Series([13.2, 14.1, 12.8, 13.5, 14.0, 13.8, 12.5, 13.9])

cv_glucose = (glucose.std() / glucose.mean()) * 100
cv_hb = (hemoglobin.std() / hemoglobin.mean()) * 100

print(f"CV glucose: {cv_glucose:.1f}%")
print(f"CV hemoglobin: {cv_hb:.1f}%")
\end{minted}

% ============================================================
\section{Frequency tables and cross-tabulations}

\subsection{One-way frequency tables}

\begin{minted}{python}
# Load Gapminder data
url = ("https://raw.githubusercontent.com/datasets/"
       "gapminder/main/data/gapminder.csv")
gapminder = pd.read_csv(url)
recent = gapminder[gapminder["year"] == 2007]

# Create life expectancy categories
recent = recent.copy()
recent["le_cat"] = pd.cut(recent["lifeExp"],
                          bins=[0, 55, 70, 85],
                          labels=["Low (<55)", "Medium (55-70)",
                                  "High (>70)"])

# Frequency table
freq = recent["le_cat"].value_counts().sort_index()
print(freq)

# With proportions
freq_table = pd.DataFrame({
    "n": freq,
    "pct": (freq / freq.sum() * 100).round(1),
    "cum_pct": (freq.cumsum() / freq.sum() * 100).round(1)
})
print(freq_table)
\end{minted}

\subsection{Two-way cross-tabulations}

\begin{minted}{python}
# Cross-tabulate life expectancy category by continent
ct = pd.crosstab(recent["continent"], recent["le_cat"], margins=True)
print(ct)

# With row percentages (percentage within each continent)
ct_pct = pd.crosstab(recent["continent"], recent["le_cat"],
                     normalize="index").round(3) * 100
print(ct_pct)
\end{minted}

\begin{clinicalnote}
Cross-tabulations are the workhorse of descriptive epidemiology. Every ``Table~1'' in a clinical paper is essentially a cross-tabulation of patient characteristics by study group. In pandas, \texttt{pd.crosstab()} builds them in one line.
\end{clinicalnote}

% ============================================================
\section{Epidemiological measures: counting disease}

\subsection{Prevalence}

Prevalence is the \emph{proportion} of a population that has a condition at a given point in time:

\[
\text{Prevalence} = \frac{\text{Number of existing cases}}{\text{Total population}} \times 100
\]

\begin{minted}{python}
# Diabetes prevalence from a health survey
n_total = 5000
n_diabetic = 625

prevalence = n_diabetic / n_total * 100
print(f"Diabetes prevalence: {prevalence:.1f}%")

# With 95% confidence interval (normal approximation)
from scipy import stats

p = n_diabetic / n_total
se = np.sqrt(p * (1 - p) / n_total)
ci_low = (p - 1.96 * se) * 100
ci_high = (p + 1.96 * se) * 100
print(f"95% CI: [{ci_low:.1f}%, {ci_high:.1f}%]")
\end{minted}

\subsection{Incidence rate}

Incidence measures \emph{new} cases over a period of time. The incidence rate uses person-time in the denominator:

\[
\text{Incidence rate} = \frac{\text{Number of new cases}}{\text{Total person-time at risk}}
\]

\begin{minted}{python}
# Tuberculosis cohort: 200 patients followed for varying durations
new_tb_cases = 18
total_person_years = 850  # sum of follow-up time for all patients

incidence_rate = new_tb_cases / total_person_years
print(f"TB incidence rate: {incidence_rate:.4f} per person-year")
print(f"                 = {incidence_rate * 1000:.1f} per 1,000 person-years")

# 95% CI for incidence rate (Poisson approximation)
ir_se = np.sqrt(new_tb_cases) / total_person_years
ci_low = (incidence_rate - 1.96 * ir_se) * 1000
ci_high = (incidence_rate + 1.96 * ir_se) * 1000
print(f"95% CI: [{ci_low:.1f}, {ci_high:.1f}] per 1,000 person-years")
\end{minted}

\begin{warningbox}
Prevalence and incidence answer different questions. Prevalence tells you the \emph{burden} of disease (how many people are sick right now). Incidence tells you the \emph{risk} of getting sick. A disease can have low incidence but high prevalence if patients survive a long time (e.g., Type~2 diabetes). A disease with high incidence and low prevalence kills quickly (e.g., Ebola).
\end{warningbox}

\subsection{Mortality rate}

\begin{minted}{python}
# Crude mortality rate
deaths = 1250
population = 500_000
mortality_rate = deaths / population * 100_000
print(f"Crude mortality rate: {mortality_rate:.1f} per 100,000")

# Case fatality rate (CFR)
total_cases = 3200
deaths_among_cases = 128
cfr = deaths_among_cases / total_cases * 100
print(f"Case fatality rate: {cfr:.1f}%")
\end{minted}

\begin{clinicalnote}
The case fatality rate (CFR) is often misinterpreted during epidemics. Early in an outbreak, the denominator (total cases) is underestimated because many mild cases are not tested. This makes the CFR appear higher than the true infection fatality rate (IFR). During early COVID-19, reported CFR was $\sim$3--4\%, but seroprevalence studies later estimated the IFR at $\sim$0.5--1\%.
\end{clinicalnote}

% ============================================================
\section{Measures of association}

\subsection{Risk ratio (relative risk)}

The risk ratio compares the probability of disease between exposed and unexposed groups:

\[
\text{RR} = \frac{\text{Risk in exposed}}{\text{Risk in unexposed}} = \frac{a/(a+b)}{c/(c+d)}
\]

\begin{minted}{python}
# Smoking and lung cancer (hypothetical cohort)
# Exposed (smokers):     a=80 cancer, b=920 no cancer
# Unexposed (non-smokers): c=10 cancer, d=990 no cancer
a, b, c, d = 80, 920, 10, 990

risk_exposed = a / (a + b)
risk_unexposed = c / (c + d)
rr = risk_exposed / risk_unexposed

print(f"Risk in smokers:     {risk_exposed:.3f} ({risk_exposed*100:.1f}%)")
print(f"Risk in non-smokers: {risk_unexposed:.3f} ({risk_unexposed*100:.1f}%)")
print(f"Risk Ratio: {rr:.2f}")

# 95% CI for log(RR)
import math
se_ln_rr = math.sqrt(1/a - 1/(a+b) + 1/c - 1/(c+d))
ln_rr = math.log(rr)
ci_low = math.exp(ln_rr - 1.96 * se_ln_rr)
ci_high = math.exp(ln_rr + 1.96 * se_ln_rr)
print(f"95% CI: [{ci_low:.2f}, {ci_high:.2f}]")
\end{minted}

\subsection{Odds ratio}

The odds ratio is used in case-control studies where you cannot directly compute risk:

\[
\text{OR} = \frac{a \times d}{b \times c}
\]

\begin{minted}{python}
# Case-control study: obesity and knee osteoarthritis
# Cases (OA):      a=120 obese, c=80 non-obese
# Controls (no OA): b=60 obese,  d=140 non-obese
a, b, c, d = 120, 60, 80, 140

odds_ratio = (a * d) / (b * c)
print(f"Odds Ratio: {odds_ratio:.2f}")

# 95% CI
se_ln_or = math.sqrt(1/a + 1/b + 1/c + 1/d)
ln_or = math.log(odds_ratio)
ci_low = math.exp(ln_or - 1.96 * se_ln_or)
ci_high = math.exp(ln_or + 1.96 * se_ln_or)
print(f"95% CI: [{ci_low:.2f}, {ci_high:.2f}]")
\end{minted}

\begin{clinicalnote}
The odds ratio approximates the risk ratio when the disease is rare ($<10\%$ prevalence in both groups). For common outcomes, the OR exaggerates the association. An OR of 2.0 for a condition with 40\% baseline prevalence does not mean ``twice the risk'' --- the actual RR would be lower.
\end{clinicalnote}

\subsection{Reusable functions}

\begin{minted}{python}
def risk_ratio(a, b, c, d, alpha=0.05):
    """Compute risk ratio with confidence interval.

    a = exposed with disease
    b = exposed without disease
    c = unexposed with disease
    d = unexposed without disease
    """
    from scipy.stats import norm
    z = norm.ppf(1 - alpha / 2)

    rr = (a / (a + b)) / (c / (c + d))
    se = math.sqrt(1/a - 1/(a+b) + 1/c - 1/(c+d))
    ci = (math.exp(math.log(rr) - z * se),
          math.exp(math.log(rr) + z * se))

    return {"RR": round(rr, 3), "CI_low": round(ci[0], 3),
            "CI_high": round(ci[1], 3)}


def odds_ratio(a, b, c, d, alpha=0.05):
    """Compute odds ratio with confidence interval."""
    from scipy.stats import norm
    z = norm.ppf(1 - alpha / 2)

    o_r = (a * d) / (b * c)
    se = math.sqrt(1/a + 1/b + 1/c + 1/d)
    ci = (math.exp(math.log(o_r) - z * se),
          math.exp(math.log(o_r) + z * se))

    return {"OR": round(o_r, 3), "CI_low": round(ci[0], 3),
            "CI_high": round(ci[1], 3)}


# Example
print(risk_ratio(80, 920, 10, 990))
print(odds_ratio(120, 60, 80, 140))
\end{minted}

% ============================================================
\section{Age-standardized rates}

Crude rates can be misleading when comparing populations with different age structures. A country with an older population will have a higher crude mortality rate even if it is ``healthier'' at every age. Age standardization removes this confounding.

\subsection{Direct standardization}

\begin{minted}{python}
# Compare crude vs age-standardized mortality between two regions
# Region A (young population), Region B (old population)

data_a = pd.DataFrame({
    "age_group": ["0-14", "15-44", "45-64", "65+"],
    "population": [50000, 80000, 40000, 10000],
    "deaths": [50, 80, 200, 300]
})
data_a["rate"] = data_a["deaths"] / data_a["population"] * 100000

data_b = pd.DataFrame({
    "age_group": ["0-14", "15-44", "45-64", "65+"],
    "population": [15000, 40000, 50000, 45000],
    "deaths": [20, 50, 280, 1500]
})
data_b["rate"] = data_b["deaths"] / data_b["population"] * 100000

# Standard population (e.g., WHO World Standard)
standard_pop = pd.DataFrame({
    "age_group": ["0-14", "15-44", "45-64", "65+"],
    "std_weight": [0.26, 0.42, 0.22, 0.10]
})

# Crude rates
crude_a = data_a["deaths"].sum() / data_a["population"].sum() * 100000
crude_b = data_b["deaths"].sum() / data_b["population"].sum() * 100000
print(f"Crude rate Region A: {crude_a:.1f} per 100,000")
print(f"Crude rate Region B: {crude_b:.1f} per 100,000")

# Age-standardized rates
asr_a = (data_a["rate"] * standard_pop["std_weight"]).sum()
asr_b = (data_b["rate"] * standard_pop["std_weight"]).sum()
print(f"Age-standardized rate Region A: {asr_a:.1f} per 100,000")
print(f"Age-standardized rate Region B: {asr_b:.1f} per 100,000")
\end{minted}

\begin{warningbox}
Always report which standard population you used (WHO World Standard, European Standard, US 2000 Standard). Different standards give different results. If you compare your rates with published rates, you must use the same standard population.
\end{warningbox}

% ============================================================
\section{Practical: computing descriptive statistics with real data}

\subsection{Descriptive statistics on Gapminder data}

\begin{minted}{python}
import pandas as pd
import numpy as np
from scipy import stats

# Load Gapminder
url = ("https://raw.githubusercontent.com/datasets/"
       "gapminder/main/data/gapminder.csv")
df = pd.read_csv(url)

# Summary statistics by continent (2007)
recent = df[df["year"] == 2007].copy()
summary = recent.groupby("continent")["lifeExp"].agg(
    ["count", "mean", "median", "std",
     lambda x: x.quantile(0.25),
     lambda x: x.quantile(0.75)]
)
summary.columns = ["n", "mean", "median", "sd", "Q1", "Q3"]
print(summary.round(1))
\end{minted}

\subsection{Computing malaria prevalence by region}

\begin{minted}{python}
# Malaria estimated cases from Our World in Data
url = ("https://raw.githubusercontent.com/owid/"
       "owid-datasets/master/datasets/"
       "Malaria%20incidence%20-%20IHME/Malaria%20incidence%20-%20IHME.csv")
# Alternative: use WHO GHO data directly
# url = "https://apps.who.int/gho/athena/api/GHO/MALARIA_EST_CASES?format=csv"

# If the above URL does not work, use Our World in Data's GitHub
url2 = ("https://raw.githubusercontent.com/owid/etl/master/etl/steps/"
        "data/garden/who/2024-07-30/gho.py")  # metadata only

# For a reliable source, use the preprocessed OWID malaria dataset
owid_url = ("https://raw.githubusercontent.com/owid/"
            "owid-datasets/master/datasets/"
            "Estimated%20malaria%20incidence%20rate%20"
            "(per%201%2C000%20population%20at%20risk)%20-%20WHO/"
            "Estimated%20malaria%20incidence%20rate%20"
            "(per%201%2C000%20population%20at%20risk)%20-%20WHO.csv")

# Simulated malaria data for practice
np.random.seed(42)
countries = (["Nigeria", "DRC", "Mozambique", "Uganda", "Ghana"] * 3 +
             ["India", "Indonesia", "Myanmar", "Pakistan", "Ethiopia"] * 3)
regions = (["West Africa"] * 5 + ["Central Africa"] * 5 +
           ["East Africa"] * 5 +
           ["South Asia"] * 5 + ["Southeast Asia"] * 5 +
           ["East Africa"] * 5)

malaria = pd.DataFrame({
    "country": countries,
    "region": ["West Africa", "Central Africa", "East Africa",
               "East Africa", "West Africa"] * 6,
    "year": [2018]*5 + [2019]*5 + [2020]*5 +
            [2018]*5 + [2019]*5 + [2020]*5,
    "population_at_risk": np.random.randint(10_000_000, 200_000_000, 30),
    "estimated_cases": np.random.randint(500_000, 50_000_000, 30)
})

# Compute prevalence per 1,000 population at risk
malaria["prevalence_per_1000"] = (
    malaria["estimated_cases"] / malaria["population_at_risk"] * 1000
)

# Summary by region
region_summary = malaria.groupby("region").agg(
    total_cases=("estimated_cases", "sum"),
    total_pop=("population_at_risk", "sum"),
    n_countries=("country", "nunique")
)
region_summary["prevalence_per_1000"] = (
    region_summary["total_cases"] / region_summary["total_pop"] * 1000
)
print(region_summary.round(1))

# Trend by year
year_trend = malaria.groupby("year").agg(
    total_cases=("estimated_cases", "sum"),
    total_pop=("population_at_risk", "sum")
)
year_trend["prevalence_per_1000"] = (
    year_trend["total_cases"] / year_trend["total_pop"] * 1000
)
print(year_trend.round(1))
\end{minted}

\begin{exercise}
Using the Gapminder dataset and the epidemiological functions defined above:
\begin{enumerate}
  \item Compute summary statistics (mean, median, SD, IQR) for life expectancy, GDP per capita, and population for each continent in 2007.
  \item Create a cross-tabulation of continent vs.\ life expectancy category (Low/Medium/High). Compute row percentages.
  \item A cohort of 10\,000 people was followed for 5 years. During this time, 450 developed Type~2 diabetes. Compute the cumulative incidence and the incidence rate (assuming an average of 4.5 years of follow-up per person).
  \item In a case-control study of 200 cases and 200 controls, 140 cases and 80 controls were exposed to a risk factor. Compute the odds ratio and its 95\% CI. Interpret the result.
  \item Country~A has a crude mortality rate of 350 per 100\,000 and Country~B has 580 per 100\,000. Using the age-specific rates and standard population provided above, compute age-standardized rates. Does the ranking change after standardization?
  \item Using Our World in Data or WHO data, compute the malaria incidence rate for 5~African countries over 3 years. Which country has the highest burden? Is the trend increasing or decreasing?
  \item Write a function \texttt{prevalence\_ci(cases, total, confidence=0.95)} that returns the prevalence and its confidence interval. Test it with: 250 diabetic patients out of 2\,000 screened.
\end{enumerate}
\end{exercise}

% ============================================================
\section{Chapter summary}

\begin{itemize}
  \item Use the \textbf{median} for skewed distributions (lab values, costs, length of stay) and the \textbf{mean} for symmetric distributions.
  \item The \textbf{IQR} is a robust measure of spread; the \textbf{SD} is sensitive to outliers.
  \item \textbf{Prevalence} = existing cases / population. \textbf{Incidence rate} = new cases / person-time.
  \item The \textbf{risk ratio} compares risks in cohort studies; the \textbf{odds ratio} is used in case-control studies.
  \item The OR approximates the RR only when the outcome is rare ($<10\%$).
  \item \textbf{Age standardization} removes confounding by age when comparing populations.
  \item Always report confidence intervals alongside point estimates.
  \item \texttt{pd.crosstab()} builds the two-way tables that form the basis of epidemiological analysis.
\end{itemize}

\chapter{Visualization for health data}
\label{ch:visualization}

\begin{quote}
\emph{``The greatest value of a picture is when it forces us to notice what we never expected to see.'' --- John Tukey}
\end{quote}

% ============================================================
\section{Principles of health data visualization}

A good health figure answers a question. A bad one decorates a slide. Before writing any plotting code, ask yourself:
\begin{enumerate}
  \item \textbf{What is the message?} ``Malaria incidence is declining in Sub-Saharan Africa'' --- not ``here is some data.''
  \item \textbf{Who is the audience?} Clinicians need different plots than policymakers.
  \item \textbf{What comparison matters?} Between groups? Over time? Against a threshold?
\end{enumerate}

Rules that apply to every health figure:
\begin{itemize}
  \item Label axes with units (``Systolic BP (mmHg)'', not ``BP'').
  \item Use colorblind-friendly palettes.
  \item Do not use 3D charts, pie charts, or dual y-axes unless you have a very good reason.
  \item Include sample sizes in titles or annotations.
  \item Start y-axes at zero for bar charts. For line plots, zero is optional if the range is narrow.
\end{itemize}

\begin{clinicalnote}
In clinical publications, figures are often the first (and sometimes the only) thing reviewers examine closely. A clear, well-labeled figure can carry a paper. A misleading one can sink it. The Lancet and BMJ both have specific figure guidelines --- study them before submitting.
\end{clinicalnote}

% ============================================================
\section{Setup: matplotlib and seaborn}

\begin{minted}{python}
import pandas as pd
import numpy as np
import matplotlib.pyplot as plt
import seaborn as sns

# Set a clean style for all plots
sns.set_theme(style="whitegrid", font_scale=1.1)
plt.rcParams["figure.figsize"] = (8, 5)
plt.rcParams["figure.dpi"] = 100

# Colorblind-friendly palette
cb_palette = ["#0072B2", "#D55E00", "#009E73",
              "#CC79A7", "#F0E442", "#56B4E9"]
sns.set_palette(cb_palette)
\end{minted}

\begin{pythontip}
The six-color palette above is from Bang Wong's ``Points of view: Color blindness'' (Nature Methods, 2011). It is distinguishable by people with the most common forms of color vision deficiency. Use it as your default.
\end{pythontip}

% ============================================================
\section{Histograms and density plots}

Histograms show the \emph{distribution} of a single variable. In health, this tells you whether a measurement follows a normal distribution, is skewed, or is bimodal.

\begin{minted}{python}
# Load Gapminder data
url = ("https://raw.githubusercontent.com/datasets/"
       "gapminder/main/data/gapminder.csv")
gapminder = pd.read_csv(url)
recent = gapminder[gapminder["year"] == 2007]

# Histogram of life expectancy
fig, axes = plt.subplots(1, 2, figsize=(12, 5))

# Left: histogram
axes[0].hist(recent["lifeExp"], bins=20, edgecolor="white", alpha=0.8)
axes[0].set_xlabel("Life expectancy (years)")
axes[0].set_ylabel("Number of countries")
axes[0].set_title("Distribution of life expectancy, 2007")
axes[0].axvline(recent["lifeExp"].median(), color="red",
                linestyle="--", label=f'Median: {recent["lifeExp"].median():.1f}')
axes[0].legend()

# Right: kernel density estimate (smooth version)
sns.kdeplot(data=recent, x="lifeExp", hue="continent",
            fill=True, alpha=0.3, ax=axes[1])
axes[1].set_xlabel("Life expectancy (years)")
axes[1].set_title("Life expectancy by continent, 2007")

plt.tight_layout()
plt.show()
\end{minted}

\begin{minted}{python}
# Clinical example: distribution of fasting glucose
np.random.seed(42)
normal_glucose = np.random.normal(95, 10, 800)
prediabetic = np.random.normal(115, 8, 150)
diabetic = np.random.normal(180, 40, 50)
all_glucose = np.concatenate([normal_glucose, prediabetic, diabetic])

fig, ax = plt.subplots(figsize=(10, 5))
ax.hist(all_glucose, bins=40, edgecolor="white", alpha=0.7)
ax.axvline(100, color="orange", linestyle="--", linewidth=2,
           label="Normal threshold (100 mg/dL)")
ax.axvline(126, color="red", linestyle="--", linewidth=2,
           label="Diabetic threshold (126 mg/dL)")
ax.set_xlabel("Fasting glucose (mg/dL)")
ax.set_ylabel("Number of patients")
ax.set_title("Distribution of fasting glucose (n=1,000)")
ax.legend()
plt.tight_layout()
plt.show()
\end{minted}

% ============================================================
\section{Box plots}

Box plots compare distributions across groups. The box shows the IQR (25th--75th percentile), the line is the median, and the whiskers extend to 1.5$\times$IQR. Points beyond the whiskers are outliers.

\begin{minted}{python}
# Blood pressure by continent
fig, ax = plt.subplots(figsize=(10, 6))
sns.boxplot(data=recent, x="continent", y="lifeExp", ax=ax)
ax.set_xlabel("Continent")
ax.set_ylabel("Life expectancy (years)")
ax.set_title("Life expectancy by continent, 2007 (n={})".format(len(recent)))
plt.tight_layout()
plt.show()
\end{minted}

\begin{minted}{python}
# Clinical example: BMI by sex and age group
np.random.seed(42)
n = 500
clinical = pd.DataFrame({
    "sex": np.random.choice(["Male", "Female"], n),
    "age_group": np.random.choice(["18-39", "40-59", "60+"], n,
                                  p=[0.3, 0.4, 0.3]),
    "bmi": np.random.normal(27, 5, n).clip(15, 50)
})

fig, ax = plt.subplots(figsize=(10, 6))
sns.boxplot(data=clinical, x="age_group", y="bmi", hue="sex",
            order=["18-39", "40-59", "60+"], ax=ax)
ax.axhline(25, color="orange", linestyle="--", alpha=0.7,
           label="Overweight threshold")
ax.axhline(30, color="red", linestyle="--", alpha=0.7,
           label="Obese threshold")
ax.set_xlabel("Age group")
ax.set_ylabel("BMI (kg/m²)")
ax.set_title(f"BMI by sex and age group (n={n})")
ax.legend()
plt.tight_layout()
plt.show()
\end{minted}

\begin{warningbox}
Box plots hide the shape of the distribution. A bimodal distribution looks the same as a unimodal one in a box plot. Consider using a \textbf{violin plot} (\texttt{sns.violinplot}) or a \textbf{strip plot} (\texttt{sns.stripplot}) overlaid on the box plot for small samples.
\end{warningbox}

% ============================================================
\section{Bar charts}

Bar charts compare \emph{counts} or \emph{rates} across categories. They are the standard choice for disease prevalence comparisons.

\begin{minted}{python}
# Life expectancy comparison: bottom 15 countries (2007)
bottom15 = recent.nsmallest(15, "lifeExp")

fig, ax = plt.subplots(figsize=(10, 7))
ax.barh(bottom15["country"], bottom15["lifeExp"], color="#D55E00")
ax.set_xlabel("Life expectancy (years)")
ax.set_title("15 countries with lowest life expectancy, 2007")
ax.invert_yaxis()  # highest at top
for i, v in enumerate(bottom15["lifeExp"]):
    ax.text(v + 0.5, i, f"{v:.1f}", va="center", fontsize=9)
plt.tight_layout()
plt.show()
\end{minted}

\begin{minted}{python}
# Disease prevalence comparison by country (simulated)
diseases = pd.DataFrame({
    "country": ["Nigeria", "DRC", "Tanzania", "Kenya", "Ghana",
                "South Africa", "Ethiopia", "Uganda"],
    "malaria_prev": [27.3, 31.2, 7.5, 5.8, 15.2, 0.5, 1.2, 9.7],
    "hiv_prev": [1.3, 0.7, 4.7, 4.9, 1.7, 19.0, 1.0, 5.4],
})

fig, ax = plt.subplots(figsize=(10, 6))
x = np.arange(len(diseases))
width = 0.35
ax.bar(x - width/2, diseases["malaria_prev"], width,
       label="Malaria", color="#0072B2")
ax.bar(x + width/2, diseases["hiv_prev"], width,
       label="HIV", color="#D55E00")
ax.set_xticks(x)
ax.set_xticklabels(diseases["country"], rotation=45, ha="right")
ax.set_ylabel("Prevalence (%)")
ax.set_title("Malaria vs HIV prevalence by country")
ax.legend()
plt.tight_layout()
plt.show()
\end{minted}

% ============================================================
\section{Line plots: time trends and epidemic curves}

Line plots show how a measurement changes over time. They are essential for tracking disease trends and epidemic progression.

\subsection{Time trends}

\begin{minted}{python}
# Life expectancy trends by continent
continent_trend = gapminder.groupby(
    ["year", "continent"])["lifeExp"].mean().reset_index()

fig, ax = plt.subplots(figsize=(10, 6))
for continent in continent_trend["continent"].unique():
    subset = continent_trend[continent_trend["continent"] == continent]
    ax.plot(subset["year"], subset["lifeExp"],
            marker="o", markersize=4, label=continent)

ax.set_xlabel("Year")
ax.set_ylabel("Mean life expectancy (years)")
ax.set_title("Life expectancy trends by continent, 1952-2007")
ax.legend(title="Continent")
plt.tight_layout()
plt.show()
\end{minted}

\subsection{Epidemic curves}

\begin{minted}{python}
# Simulate an epidemic curve (SIR-like shape)
np.random.seed(42)
days = pd.date_range("2023-01-01", periods=120, freq="D")
# Simulate daily new cases with an SIR-like peak
t = np.arange(120)
peak = 45
cases = np.maximum(0, 200 * np.exp(-0.5 * ((t - peak) / 15) ** 2)
                   + np.random.normal(0, 10, 120)).astype(int)

epi = pd.DataFrame({"date": days, "new_cases": cases})
epi["rolling_7d"] = epi["new_cases"].rolling(7).mean()
epi["cumulative"] = epi["new_cases"].cumsum()

fig, axes = plt.subplots(2, 1, figsize=(12, 8), sharex=True)

# Top: daily cases + 7-day average
axes[0].bar(epi["date"], epi["new_cases"], alpha=0.4,
            color="#0072B2", label="Daily cases")
axes[0].plot(epi["date"], epi["rolling_7d"], color="#D55E00",
             linewidth=2, label="7-day average")
axes[0].set_ylabel("New cases per day")
axes[0].set_title("Epidemic curve: daily new cases")
axes[0].legend()

# Bottom: cumulative cases
axes[1].plot(epi["date"], epi["cumulative"], color="#009E73",
             linewidth=2)
axes[1].set_xlabel("Date")
axes[1].set_ylabel("Cumulative cases")
axes[1].set_title("Cumulative case count")

plt.tight_layout()
plt.show()
\end{minted}

\begin{clinicalnote}
During outbreaks, the epidemic curve (``epi curve'') is the single most important visualization. Its shape tells you the transmission pattern: a point-source outbreak has a sharp peak; a propagated outbreak has successive waves. The 7-day rolling average smooths out reporting delays and weekend effects.
\end{clinicalnote}

% ============================================================
\section{Scatter plots}

Scatter plots show the relationship between two continuous variables.

\begin{minted}{python}
# GDP per capita vs life expectancy (the classic Gapminder plot)
fig, ax = plt.subplots(figsize=(10, 7))
for continent in recent["continent"].unique():
    subset = recent[recent["continent"] == continent]
    ax.scatter(subset["gdpPercap"], subset["lifeExp"],
               s=subset["pop"] / 1e6,  # bubble size = population
               alpha=0.6, label=continent)

ax.set_xscale("log")
ax.set_xlabel("GDP per capita (log scale, USD)")
ax.set_ylabel("Life expectancy (years)")
ax.set_title("Wealth vs Health, 2007 (bubble size = population)")
ax.legend(title="Continent")
plt.tight_layout()
plt.show()
\end{minted}

\begin{minted}{python}
# Clinical scatter: BMI vs systolic BP with regression line
np.random.seed(42)
n = 200
bmi = np.random.normal(28, 5, n).clip(18, 45)
sbp = 80 + 1.8 * bmi + np.random.normal(0, 12, n)

fig, ax = plt.subplots(figsize=(8, 6))
ax.scatter(bmi, sbp, alpha=0.5, s=30, color="#0072B2")

# Add regression line
from numpy.polynomial.polynomial import polyfit
b, m = polyfit(bmi, sbp, 1)
x_line = np.linspace(bmi.min(), bmi.max(), 100)
ax.plot(x_line, b + m * x_line, color="#D55E00", linewidth=2,
        label=f"y = {m:.1f}x + {b:.1f}")

# Correlation coefficient
from scipy.stats import pearsonr
r, p = pearsonr(bmi, sbp)
ax.set_xlabel("BMI (kg/m²)")
ax.set_ylabel("Systolic BP (mmHg)")
ax.set_title(f"BMI vs Systolic BP (r={r:.2f}, p<0.001, n={n})")
ax.legend()
plt.tight_layout()
plt.show()
\end{minted}

\begin{pythontip}
Seaborn's \texttt{sns.regplot()} and \texttt{sns.lmplot()} add regression lines and confidence bands automatically. Use \texttt{sns.lmplot()} when you want to split by a categorical variable (e.g., regression by sex).
\end{pythontip}

% ============================================================
\section{Kaplan-Meier survival curves}

Survival analysis tracks time until an event (death, relapse, infection). The Kaplan-Meier curve estimates the probability of surviving beyond each time point.

\begin{minted}{python}
# Install lifelines if needed: !pip install lifelines
from lifelines import KaplanMeierFitter

# Simulated survival data: two treatment groups
np.random.seed(42)
n_per_group = 100

# Treatment group: longer survival
treatment_time = np.random.exponential(24, n_per_group).clip(0, 60)
treatment_event = np.random.binomial(1, 0.7, n_per_group)

# Control group: shorter survival
control_time = np.random.exponential(15, n_per_group).clip(0, 60)
control_event = np.random.binomial(1, 0.8, n_per_group)

fig, ax = plt.subplots(figsize=(10, 6))

# Fit and plot treatment group
kmf_treatment = KaplanMeierFitter()
kmf_treatment.fit(treatment_time, event_observed=treatment_event,
                  label="Treatment (n=100)")
kmf_treatment.plot_survival_function(ax=ax, ci_show=True)

# Fit and plot control group
kmf_control = KaplanMeierFitter()
kmf_control.fit(control_time, event_observed=control_event,
                label="Control (n=100)")
kmf_control.plot_survival_function(ax=ax, ci_show=True)

ax.set_xlabel("Time (months)")
ax.set_ylabel("Survival probability")
ax.set_title("Kaplan-Meier survival curves by treatment group")
ax.set_ylim(0, 1.05)

# Add median survival lines
for kmf, color in [(kmf_treatment, "#0072B2"), (kmf_control, "#D55E00")]:
    median = kmf.median_survival_time_
    if not np.isinf(median):
        ax.axhline(0.5, color="gray", linestyle=":", alpha=0.5)
        ax.axvline(median, color=color, linestyle=":", alpha=0.5)

plt.tight_layout()
plt.show()

# Print median survival
print(f"Median survival (treatment): "
      f"{kmf_treatment.median_survival_time_:.1f} months")
print(f"Median survival (control): "
      f"{kmf_control.median_survival_time_:.1f} months")
\end{minted}

\begin{warningbox}
If \texttt{lifelines} is not installed, run \texttt{!pip install lifelines} in your notebook. It is not pre-installed in Google Colab.
\end{warningbox}

\begin{clinicalnote}
Kaplan-Meier curves handle \textbf{censored} observations --- patients who are still alive at the end of the study or who were lost to follow-up. Without proper censoring, you would underestimate survival. The vertical tick marks on the curve show censored observations.
\end{clinicalnote}

% ============================================================
\section{Heatmaps}

Heatmaps display matrices of values using color intensity. Two common uses in health:

\subsection{Correlation matrices}

\begin{minted}{python}
# Correlation between health indicators
np.random.seed(42)
n = 300
health = pd.DataFrame({
    "BMI": np.random.normal(27, 5, n),
    "Systolic_BP": np.random.normal(130, 18, n),
    "Diastolic_BP": np.random.normal(82, 10, n),
    "Glucose": np.random.normal(105, 25, n),
    "Cholesterol": np.random.normal(200, 40, n),
    "HbA1c": np.random.normal(5.8, 1.2, n),
    "Heart_Rate": np.random.normal(72, 12, n),
    "Age": np.random.normal(55, 15, n).clip(18, 90)
})

# Add realistic correlations
health["Systolic_BP"] += 0.8 * health["BMI"]
health["Glucose"] += 15 * health["HbA1c"]
health["Diastolic_BP"] += 0.4 * health["Systolic_BP"]
health["Cholesterol"] += 0.5 * health["BMI"]

corr = health.corr()

fig, ax = plt.subplots(figsize=(9, 8))
mask = np.triu(np.ones_like(corr, dtype=bool))  # hide upper triangle
sns.heatmap(corr, mask=mask, annot=True, fmt=".2f",
            cmap="RdBu_r", center=0, vmin=-1, vmax=1,
            square=True, ax=ax)
ax.set_title("Correlation matrix of health indicators")
plt.tight_layout()
plt.show()
\end{minted}

\subsection{Disease co-occurrence}

\begin{minted}{python}
# Co-occurrence of chronic conditions
conditions = ["Hypertension", "Diabetes", "COPD",
              "Heart Failure", "CKD", "Obesity"]
np.random.seed(42)
cooccurrence = np.random.randint(5, 60, (6, 6))
cooccurrence = (cooccurrence + cooccurrence.T) // 2  # make symmetric
np.fill_diagonal(cooccurrence, [320, 180, 95, 75, 110, 210])

co_df = pd.DataFrame(cooccurrence,
                     index=conditions, columns=conditions)

fig, ax = plt.subplots(figsize=(9, 8))
sns.heatmap(co_df, annot=True, fmt="d", cmap="YlOrRd",
            square=True, ax=ax)
ax.set_title("Disease co-occurrence matrix\n"
             "(diagonal = total cases, off-diagonal = co-occurring)")
plt.tight_layout()
plt.show()
\end{minted}

% ============================================================
\section{Practical: build a 4-panel health dashboard}

Combining multiple plots into a single figure is essential for reports and publications. Here is a complete 4-panel dashboard using real data:

\begin{minted}{python}
import pandas as pd
import numpy as np
import matplotlib.pyplot as plt
import seaborn as sns

sns.set_theme(style="whitegrid", font_scale=1.0)

# Load data
url = ("https://raw.githubusercontent.com/datasets/"
       "gapminder/main/data/gapminder.csv")
df = pd.read_csv(url)
recent = df[df["year"] == 2007].copy()

# Create the 4-panel figure
fig, axes = plt.subplots(2, 2, figsize=(14, 10))
fig.suptitle("Global Health Dashboard, 2007",
             fontsize=16, fontweight="bold", y=1.02)

# ----------------------------------------------------------
# Panel A: Life expectancy distribution by continent
# ----------------------------------------------------------
ax = axes[0, 0]
sns.boxplot(data=recent, x="continent", y="lifeExp", ax=ax,
            palette="Set2")
ax.set_xlabel("")
ax.set_ylabel("Life expectancy (years)")
ax.set_title("A. Life expectancy by continent")
ax.tick_params(axis="x", rotation=30)

# ----------------------------------------------------------
# Panel B: GDP vs Life expectancy scatter
# ----------------------------------------------------------
ax = axes[0, 1]
for continent in recent["continent"].unique():
    subset = recent[recent["continent"] == continent]
    ax.scatter(subset["gdpPercap"], subset["lifeExp"],
               s=subset["pop"] / 2e6, alpha=0.6,
               label=continent)
ax.set_xscale("log")
ax.set_xlabel("GDP per capita (USD, log scale)")
ax.set_ylabel("Life expectancy (years)")
ax.set_title("B. Wealth vs health")
ax.legend(fontsize=8, title="Continent", title_fontsize=8)

# ----------------------------------------------------------
# Panel C: Life expectancy trends over time
# ----------------------------------------------------------
ax = axes[1, 0]
trend = df.groupby(["year", "continent"])["lifeExp"].mean().reset_index()
for continent in trend["continent"].unique():
    subset = trend[trend["continent"] == continent]
    ax.plot(subset["year"], subset["lifeExp"],
            marker="o", markersize=3, label=continent)
ax.set_xlabel("Year")
ax.set_ylabel("Mean life expectancy (years)")
ax.set_title("C. Life expectancy trends, 1952-2007")
ax.legend(fontsize=8, title="Continent", title_fontsize=8)

# ----------------------------------------------------------
# Panel D: Bottom 10 countries bar chart
# ----------------------------------------------------------
ax = axes[1, 1]
bottom10 = recent.nsmallest(10, "lifeExp")
ax.barh(bottom10["country"], bottom10["lifeExp"], color="#D55E00")
ax.set_xlabel("Life expectancy (years)")
ax.set_title("D. 10 lowest life expectancy countries")
ax.invert_yaxis()
for i, v in enumerate(bottom10["lifeExp"]):
    ax.text(v + 0.3, i, f"{v:.1f}", va="center", fontsize=8)

plt.tight_layout()
plt.savefig("health_dashboard_2007.png", dpi=150, bbox_inches="tight")
plt.show()
print("Dashboard saved as health_dashboard_2007.png")
\end{minted}

\begin{exercise}
\begin{enumerate}
  \item Modify the dashboard to show 2002 data instead of 2007. Does the story change?
  \item Add a 5th panel showing a histogram of GDP per capita (log-transformed) with a vertical line at the world median. Use \texttt{fig, axes = plt.subplots(2, 3)} or \texttt{plt.subplot2grid()}.
  \item Create a clinical dashboard for a simulated cohort of 500 patients with the following panels:
    \begin{itemize}
      \item Histogram of BMI with overweight/obese thresholds
      \item Box plot of systolic BP by sex
      \item Scatter plot of age vs.\ HbA1c with a diabetic threshold line at 6.5\%
      \item Bar chart of diagnosis frequency
    \end{itemize}
  \item Using Our World in Data (\url{https://github.com/owid/owid-datasets}), download a dataset on malaria or tuberculosis. Build an epidemic dashboard with:
    \begin{itemize}
      \item Line plot of incidence over time for 5 countries
      \item Bar chart of the most recent year's incidence
      \item Scatter plot of incidence vs.\ GDP
      \item Heatmap of correlation between health indicators
    \end{itemize}
  \item Create a Kaplan-Meier plot comparing survival between 3 treatment groups (drug A, drug B, placebo) using simulated data. Add median survival annotations.
  \item Build a ``pre/post intervention'' visualization: show a health indicator (e.g., malaria cases) as a line plot with a vertical line marking the intervention date. Shade the pre-intervention and post-intervention periods in different colors.
\end{enumerate}
\end{exercise}

% ============================================================
\section{Chapter summary}

\begin{itemize}
  \item Always label axes with units, include sample sizes, and use colorblind-friendly palettes.
  \item \textbf{Histograms} show distributions; add clinical threshold lines to give context.
  \item \textbf{Box plots} compare groups; consider violin plots for more detail.
  \item \textbf{Bar charts} compare counts or rates; use horizontal bars for many categories.
  \item \textbf{Line plots} show trends over time; add 7-day rolling averages for epidemic curves.
  \item \textbf{Scatter plots} reveal relationships; add regression lines and report $r$ values.
  \item \textbf{Kaplan-Meier curves} are the standard for survival data; use the \texttt{lifelines} library.
  \item \textbf{Heatmaps} display correlation matrices and co-occurrence patterns.
  \item Multi-panel figures (\texttt{plt.subplots()}) combine related plots into a coherent story.
  \item Save publication-quality figures with \texttt{plt.savefig("file.png", dpi=300, bbox\_inches="tight")}.
\end{itemize}

\chapter{Hypothesis testing}
\label{ch:hypothesis-testing}

\begin{quote}
\emph{``The purpose of a statistical test is not to prove something is true. It is to measure how surprised we should be if nothing is going on.''}
\end{quote}

% ============================================================
\section{The logic of hypothesis testing}

Every clinical question can be framed as a test between two competing statements:

\begin{itemize}
  \item \textbf{Null hypothesis ($H_0$):} There is \emph{no} effect, no difference, no association. The drug does not work. The two groups have the same blood pressure.
  \item \textbf{Alternative hypothesis ($H_1$):} There \emph{is} an effect, a difference, or an association. The drug lowers cholesterol. The treatment group has lower blood pressure than the control.
\end{itemize}

We collect data, compute a test statistic, and ask: \emph{if $H_0$ were true, how likely is it that we would see data this extreme?} That probability is the \textbf{p-value}.

\begin{clinicalnote}
Clinical framing matters. A statistician writes $H_0: \mu_1 = \mu_2$. A clinician writes ``there is no difference in systolic blood pressure between the statin group and the placebo group.'' Both say the same thing --- use whichever makes the research question clear.
\end{clinicalnote}

\subsection{The steps of a hypothesis test}

\begin{enumerate}
  \item \textbf{State the hypotheses} ($H_0$ and $H_1$).
  \item \textbf{Choose the significance level} $\alpha$ (usually 0.05).
  \item \textbf{Collect data} and compute the \textbf{test statistic}.
  \item \textbf{Compute the p-value} --- the probability of observing a result at least as extreme as yours, assuming $H_0$ is true.
  \item \textbf{Decide:} if $p < \alpha$, reject $H_0$. Otherwise, fail to reject $H_0$.
\end{enumerate}

\begin{minted}{python}
import numpy as np
import pandas as pd
from scipy import stats
import matplotlib.pyplot as plt
import seaborn as sns
\end{minted}

% ============================================================
\section{p-values: what they mean and what they don't}

\subsection{What a p-value IS}

A p-value is the probability of observing data as extreme as (or more extreme than) what you collected, \emph{assuming} $H_0$ is true. A $p = 0.03$ means: if the drug truly has no effect, there is a 3\% chance of seeing results this extreme by random chance alone.

\subsection{What a p-value is NOT}

\begin{itemize}
  \item It is \textbf{not} the probability that $H_0$ is true.
  \item It is \textbf{not} the probability that the result is due to chance.
  \item It does \textbf{not} measure the size of the effect.
  \item $p = 0.049$ is not meaningfully different from $p = 0.051$.
\end{itemize}

\begin{warningbox}
A small p-value does not mean a \emph{large} or \emph{clinically important} effect. A study with 100\,000 patients can find $p < 0.001$ for a blood pressure difference of 0.5~mmHg --- statistically significant but clinically meaningless. Always report \textbf{effect size} alongside p-values.
\end{warningbox}

\subsection{Type I and Type II errors}

\begin{itemize}
  \item \textbf{Type I error (false positive):} You reject $H_0$ when it is actually true. You conclude the drug works when it does not. Controlled by $\alpha$.
  \item \textbf{Type II error (false negative):} You fail to reject $H_0$ when $H_1$ is actually true. You miss a real effect. Related to \textbf{power} ($1 - \beta$).
\end{itemize}

\begin{minted}{python}
# Visualize the logic: sampling distribution under H0
np.random.seed(42)
null_distribution = np.random.normal(loc=0, scale=1, size=10000)

fig, ax = plt.subplots(figsize=(9, 4))
ax.hist(null_distribution, bins=60, density=True, alpha=0.7, color="steelblue")
ax.axvline(1.96, color="red", linestyle="--", label="Critical value (alpha=0.05)")
ax.axvline(-1.96, color="red", linestyle="--")
ax.fill_betweenx([0, 0.45], 1.96, 3.5, alpha=0.3, color="red", label="Rejection region")
ax.fill_betweenx([0, 0.45], -3.5, -1.96, alpha=0.3, color="red")
ax.set_xlabel("Test statistic (z)")
ax.set_ylabel("Density")
ax.set_title("Two-tailed test: rejection regions under H0")
ax.legend()
plt.tight_layout()
plt.show()
\end{minted}

% ============================================================
\section{One-sample t-test}

\textbf{Question:} Is the mean total cholesterol in our clinic patients different from the national average of 200~mg/dL?

\begin{minted}{python}
# Simulated clinic cholesterol data (mg/dL)
np.random.seed(42)
cholesterol = np.random.normal(loc=212, scale=35, size=80)

# One-sample t-test: is the mean different from 200?
t_stat, p_value = stats.ttest_1samp(cholesterol, popmean=200)
print(f"Sample mean: {cholesterol.mean():.1f} mg/dL")
print(f"t-statistic: {t_stat:.3f}")
print(f"p-value: {p_value:.4f}")

if p_value < 0.05:
    print("Reject H0: mean cholesterol differs from 200 mg/dL")
else:
    print("Fail to reject H0: no evidence mean differs from 200 mg/dL")
\end{minted}

\begin{clinicalnote}
The one-sample t-test compares a sample mean to a known reference value. In clinical practice, the reference might be a national average, a guideline threshold, or a pre-treatment baseline.
\end{clinicalnote}

% ============================================================
\section{Two-sample t-test}

\textbf{Question:} Is systolic blood pressure different between the treatment group and the control group?

\begin{minted}{python}
# Simulated clinical trial data
np.random.seed(42)
treatment = np.random.normal(loc=128, scale=15, size=60)
control = np.random.normal(loc=138, scale=16, size=55)

# Independent two-sample t-test
t_stat, p_value = stats.ttest_ind(treatment, control)
print(f"Treatment mean: {treatment.mean():.1f} mmHg")
print(f"Control mean:   {control.mean():.1f} mmHg")
print(f"Difference:     {control.mean() - treatment.mean():.1f} mmHg")
print(f"t-statistic:    {t_stat:.3f}")
print(f"p-value:        {p_value:.4f}")
\end{minted}

\subsection{Checking assumptions}

The t-test assumes approximate normality and equal variances. Check both:

\begin{minted}{python}
# Normality: Shapiro-Wilk test (H0: data is normally distributed)
_, p_treat = stats.shapiro(treatment)
_, p_ctrl = stats.shapiro(control)
print(f"Shapiro-Wilk p (treatment): {p_treat:.4f}")
print(f"Shapiro-Wilk p (control):   {p_ctrl:.4f}")

# Equal variances: Levene's test
_, p_levene = stats.levene(treatment, control)
print(f"Levene's test p: {p_levene:.4f}")

# If variances are unequal, use Welch's t-test:
t_stat, p_value = stats.ttest_ind(treatment, control, equal_var=False)
print(f"Welch's t-test p-value: {p_value:.4f}")
\end{minted}

\begin{pythontip}
Set \texttt{equal\_var=False} in \texttt{ttest\_ind()} to use Welch's t-test. It does not assume equal variances and is often the safer default.
\end{pythontip}

% ============================================================
\section{Paired t-test}

\textbf{Question:} Does a 12-week exercise intervention reduce systolic blood pressure?

Each patient is measured \emph{before} and \emph{after}. The paired t-test uses the within-patient differences.

\begin{minted}{python}
# Simulated before/after intervention data
np.random.seed(42)
n_patients = 45
bp_before = np.random.normal(loc=142, scale=12, size=n_patients)
# After intervention: mean reduction of ~8 mmHg with individual variation
bp_after = bp_before - np.random.normal(loc=8, scale=6, size=n_patients)

# Paired t-test
t_stat, p_value = stats.ttest_rel(bp_before, bp_after)
differences = bp_before - bp_after
print(f"Mean BP before:     {bp_before.mean():.1f} mmHg")
print(f"Mean BP after:      {bp_after.mean():.1f} mmHg")
print(f"Mean reduction:     {differences.mean():.1f} mmHg")
print(f"t-statistic:        {t_stat:.3f}")
print(f"p-value:            {p_value:.6f}")
\end{minted}

\begin{minted}{python}
# Visualize before/after
fig, axes = plt.subplots(1, 2, figsize=(12, 5))

# Paired lines
for i in range(n_patients):
    axes[0].plot([0, 1], [bp_before[i], bp_after[i]],
                 color="gray", alpha=0.4)
axes[0].plot([0, 1], [bp_before.mean(), bp_after.mean()],
             color="red", linewidth=3, marker="o", markersize=8)
axes[0].set_xticks([0, 1])
axes[0].set_xticklabels(["Before", "After"])
axes[0].set_ylabel("Systolic BP (mmHg)")
axes[0].set_title("Individual patient trajectories")

# Distribution of differences
axes[1].hist(differences, bins=15, edgecolor="black", alpha=0.7, color="steelblue")
axes[1].axvline(0, color="red", linestyle="--", label="No change")
axes[1].axvline(differences.mean(), color="green", linestyle="--",
                label=f"Mean = {differences.mean():.1f}")
axes[1].set_xlabel("BP reduction (mmHg)")
axes[1].set_ylabel("Frequency")
axes[1].set_title("Distribution of BP changes")
axes[1].legend()

plt.tight_layout()
plt.show()
\end{minted}

\begin{warningbox}
A common mistake is using an \emph{independent} two-sample t-test on before/after data. This ignores the pairing and loses statistical power. If the same patients are measured twice, always use a \textbf{paired} t-test.
\end{warningbox}

% ============================================================
\section{Chi-square test of independence}

\textbf{Question:} Is diabetes prevalence associated with obesity category?

The chi-square test works with \emph{categorical} variables. It compares observed counts to what we would expect if the variables were independent.

\begin{minted}{python}
# Simulated cross-tabulation data
np.random.seed(42)
n = 500
bmi_cat = np.random.choice(["Normal", "Overweight", "Obese"],
                           size=n, p=[0.3, 0.4, 0.3])
# Diabetes probability increases with obesity
diabetes_prob = {"Normal": 0.08, "Overweight": 0.18, "Obese": 0.35}
diabetes = [np.random.binomial(1, diabetes_prob[b]) for b in bmi_cat]

df_chi = pd.DataFrame({"bmi_category": bmi_cat, "diabetes": diabetes})
df_chi["diabetes_label"] = df_chi["diabetes"].map({0: "No", 1: "Yes"})

# Contingency table
contingency = pd.crosstab(df_chi["bmi_category"], df_chi["diabetes_label"])
print(contingency)
print()

# Chi-square test
chi2, p_value, dof, expected = stats.chi2_contingency(contingency)
print(f"Chi-square statistic: {chi2:.2f}")
print(f"Degrees of freedom:   {dof}")
print(f"p-value:              {p_value:.4f}")
print(f"\nExpected frequencies (if independent):")
print(pd.DataFrame(expected, index=contingency.index,
                   columns=contingency.columns).round(1))
\end{minted}

\begin{clinicalnote}
The chi-square test requires that expected cell counts are at least 5. If you have small cells (e.g., a rare disease), use Fisher's exact test instead: \texttt{stats.fisher\_exact(table)} for $2 \times 2$ tables.
\end{clinicalnote}

% ============================================================
\section{Mann-Whitney U test}

When your data is not normally distributed --- common with hospital length-of-stay, cost data, or small samples --- use the Mann-Whitney U test. It compares the \emph{ranks} of observations rather than the means.

\begin{minted}{python}
# Hospital length of stay: surgical vs non-surgical (right-skewed data)
np.random.seed(42)
los_surgical = np.random.exponential(scale=7, size=40)
los_medical = np.random.exponential(scale=4.5, size=45)

# Mann-Whitney U test
u_stat, p_value = stats.mannwhitneyu(los_surgical, los_medical,
                                      alternative="two-sided")
print(f"Median LOS (surgical): {np.median(los_surgical):.1f} days")
print(f"Median LOS (medical):  {np.median(los_medical):.1f} days")
print(f"U statistic:           {u_stat:.0f}")
print(f"p-value:               {p_value:.4f}")
\end{minted}

\begin{minted}{python}
# Visualize the skewed distributions
fig, ax = plt.subplots(figsize=(8, 4))
ax.hist(los_surgical, bins=15, alpha=0.6, label="Surgical", color="coral")
ax.hist(los_medical, bins=15, alpha=0.6, label="Medical", color="steelblue")
ax.set_xlabel("Length of stay (days)")
ax.set_ylabel("Frequency")
ax.set_title("Hospital length of stay by department")
ax.legend()
plt.tight_layout()
plt.show()
\end{minted}

\begin{pythontip}
Rule of thumb for choosing a test:
\begin{itemize}
  \item Normal data, two groups $\rightarrow$ t-test
  \item Non-normal data, two groups $\rightarrow$ Mann-Whitney U
  \item Categorical data $\rightarrow$ chi-square
  \item Paired measurements $\rightarrow$ paired t-test (or Wilcoxon signed-rank if non-normal)
\end{itemize}
\end{pythontip}

% ============================================================
\section{Multiple testing correction}

\subsection{The problem}

If you test 20 hypotheses at $\alpha = 0.05$, you expect $20 \times 0.05 = 1$ false positive \emph{even if nothing is going on}. In genomics, you might test 20\,000 genes simultaneously. Without correction, you would get 1\,000 spurious ``discoveries.''

\begin{minted}{python}
# Simulation: 1000 tests where H0 is true for all
np.random.seed(42)
n_tests = 1000
p_values = []
for _ in range(n_tests):
    # Two random samples from the SAME distribution (H0 is true)
    a = np.random.normal(0, 1, 30)
    b = np.random.normal(0, 1, 30)
    _, p = stats.ttest_ind(a, b)
    p_values.append(p)

p_values = np.array(p_values)
false_positives = (p_values < 0.05).sum()
print(f"Tests performed: {n_tests}")
print(f"False positives (p < 0.05): {false_positives}")
print(f"False positive rate: {false_positives/n_tests:.1%}")
\end{minted}

\subsection{Bonferroni correction}

The simplest fix: divide $\alpha$ by the number of tests.

\begin{minted}{python}
from statsmodels.stats.multitest import multipletests

# Bonferroni correction
rejected_bonf, pvals_corrected_bonf, _, _ = multipletests(
    p_values, alpha=0.05, method="bonferroni"
)
print(f"Bonferroni: {rejected_bonf.sum()} significant results")
\end{minted}

\subsection{Benjamini-Hochberg (FDR)}

Less conservative: controls the \emph{false discovery rate} (the proportion of rejected hypotheses that are false positives).

\begin{minted}{python}
# Benjamini-Hochberg FDR correction
rejected_fdr, pvals_corrected_fdr, _, _ = multipletests(
    p_values, alpha=0.05, method="fdr_bh"
)
print(f"FDR (BH): {rejected_fdr.sum()} significant results")
\end{minted}

\begin{clinicalnote}
In genomics and proteomics, FDR correction is standard. When you read a paper reporting ``542 differentially expressed genes at FDR $< 0.05$,'' the authors used Benjamini-Hochberg (or a similar method) to control for the thousands of simultaneous tests.
\end{clinicalnote}

\begin{warningbox}
Bonferroni is very conservative --- it minimizes false positives but can miss real effects (high Type~II error). FDR is a good middle ground for exploratory analyses. For confirmatory clinical trials, Bonferroni is often preferred because false positives are more costly.
\end{warningbox}

% ============================================================
\section{Effect size: clinical vs statistical significance}

\subsection{Cohen's d}

Cohen's $d$ measures the size of the difference in \emph{standard deviation units}:
$$d = \frac{\bar{x}_1 - \bar{x}_2}{s_{\text{pooled}}}$$

\begin{itemize}
  \item $|d| < 0.2$: negligible effect
  \item $0.2 \leq |d| < 0.5$: small effect
  \item $0.5 \leq |d| < 0.8$: medium effect
  \item $|d| \geq 0.8$: large effect
\end{itemize}

\begin{minted}{python}
def cohens_d(group1, group2):
    """Compute Cohen's d for two independent groups."""
    n1, n2 = len(group1), len(group2)
    var1, var2 = group1.var(ddof=1), group2.var(ddof=1)
    pooled_std = np.sqrt(((n1 - 1) * var1 + (n2 - 1) * var2) / (n1 + n2 - 2))
    return (group1.mean() - group2.mean()) / pooled_std

# Using our treatment vs control BP data
d = cohens_d(treatment, control)
print(f"Cohen's d: {d:.2f}")
print(f"Interpretation: {'large' if abs(d) >= 0.8 else 'medium' if abs(d) >= 0.5 else 'small' if abs(d) >= 0.2 else 'negligible'} effect")
\end{minted}

\subsection{Why effect size matters}

\begin{minted}{python}
# Large sample + tiny effect = significant p-value
np.random.seed(42)
huge_group1 = np.random.normal(120.0, 15, size=50000)
huge_group2 = np.random.normal(120.5, 15, size=50000)  # only 0.5 mmHg difference

t, p = stats.ttest_ind(huge_group1, huge_group2)
d = cohens_d(huge_group1, huge_group2)
print(f"Mean difference: {huge_group2.mean() - huge_group1.mean():.2f} mmHg")
print(f"p-value: {p:.6f}")
print(f"Cohen's d: {d:.3f}")
print("Statistically significant? Yes. Clinically meaningful? No.")
\end{minted}

\begin{clinicalnote}
In a clinical trial report, the FDA and EMA expect both p-values \emph{and} confidence intervals for the treatment effect. A drug that lowers HbA1c by 0.01\% with $p < 0.001$ will not get approved. The effect must be \textbf{clinically meaningful}.
\end{clinicalnote}

% ============================================================
\section{Practical: hypothesis testing with Framingham-style data}

\begin{minted}{python}
# Load a Framingham-style cardiovascular dataset
url = ("https://raw.githubusercontent.com/dsrscientist/"
       "dataset-for-machine-learning/master/framingham.csv")
fram = pd.read_csv(url)
print(f"Shape: {fram.shape}")
print(f"Columns: {fram.columns.tolist()}")
fram.head()
\end{minted}

\begin{minted}{python}
# Quick exploration
fram.describe()
\end{minted}

\begin{minted}{python}
# Drop rows with missing values for simplicity
fram_clean = fram.dropna()
print(f"Clean dataset: {fram_clean.shape[0]} patients")
\end{minted}

\begin{exercise}
Using the Framingham dataset:
\begin{enumerate}
  \item \textbf{One-sample t-test.} Is the mean total cholesterol in this cohort different from the US national average of 200~mg/dL?
  \item \textbf{Two-sample t-test.} Compare systolic blood pressure between patients who developed CHD (\texttt{TenYearCHD = 1}) and those who did not. Report the p-value, mean difference, and Cohen's $d$.
  \item \textbf{Chi-square test.} Create a $2 \times 2$ contingency table of \texttt{currentSmoker} vs \texttt{TenYearCHD}. Is smoking associated with 10-year CHD risk?
  \item \textbf{Mann-Whitney U.} Compare \texttt{cigsPerDay} between CHD and non-CHD groups. Why is Mann-Whitney more appropriate than a t-test here?
  \item \textbf{Multiple testing.} Run t-tests comparing CHD vs non-CHD for all continuous variables (age, totChol, sysBP, diaBP, BMI, heartRate, glucose). Apply Bonferroni and FDR correction. Which variables remain significant after each correction?
  \item \textbf{Clinical interpretation.} For the variable with the largest effect size (Cohen's $d$), discuss whether the difference is clinically meaningful.
\end{enumerate}
\end{exercise}

% ============================================================
\section{Chapter summary}

\begin{itemize}
  \item A hypothesis test measures how surprising the data would be if $H_0$ were true.
  \item The p-value is \emph{not} the probability that $H_0$ is true.
  \item One-sample t-test: compare a sample mean to a reference value.
  \item Two-sample t-test: compare means between two independent groups.
  \item Paired t-test: compare means from the same subjects measured twice.
  \item Chi-square test: test association between two categorical variables.
  \item Mann-Whitney U: non-parametric alternative when data is not normal.
  \item Multiple testing correction (Bonferroni, FDR) prevents false discoveries.
  \item Effect size (Cohen's $d$) distinguishes statistical from clinical significance.
  \item Always report both p-values and effect sizes in clinical research.
\end{itemize}

\chapter{Regression for health outcomes}
\label{ch:regression}

\begin{quote}
\emph{``Correlation tells you that two things move together. Regression tells you by how much --- and lets you adjust for confounders.''}
\end{quote}

% ============================================================
\section{From correlation to prediction}

In Chapter~6 we asked ``is there a difference?'' Now we ask ``can we predict one variable from another?'' Regression is the workhorse of clinical research: it appears in nearly every epidemiological study, every clinical trial analysis, and every risk prediction model.

\begin{minted}{python}
import numpy as np
import pandas as pd
import matplotlib.pyplot as plt
import seaborn as sns
from scipy import stats
\end{minted}

% ============================================================
\section{Simple linear regression}

\textbf{Question:} How does systolic blood pressure change with age?

\begin{minted}{python}
# Simulated clinical data
np.random.seed(42)
n = 200
age = np.random.uniform(25, 80, n)
# True relationship: SBP = 90 + 0.7 * age + noise
sbp = 90 + 0.7 * age + np.random.normal(0, 12, n)

df = pd.DataFrame({"age": age, "systolic_bp": sbp})

# Fit with scipy
slope, intercept, r_value, p_value, std_err = stats.linregress(df["age"], df["systolic_bp"])
print(f"Intercept: {intercept:.2f}")
print(f"Slope: {slope:.2f} mmHg per year of age")
print(f"R-squared: {r_value**2:.3f}")
print(f"p-value: {p_value:.2e}")
\end{minted}

\begin{minted}{python}
# Visualize
fig, ax = plt.subplots(figsize=(8, 5))
ax.scatter(df["age"], df["systolic_bp"], alpha=0.5, s=20, color="steelblue")
x_line = np.linspace(25, 80, 100)
ax.plot(x_line, intercept + slope * x_line, color="red", linewidth=2,
        label=f"SBP = {intercept:.1f} + {slope:.2f} x Age")
ax.set_xlabel("Age (years)")
ax.set_ylabel("Systolic BP (mmHg)")
ax.set_title("Simple linear regression: SBP vs Age")
ax.legend()
plt.tight_layout()
plt.show()
\end{minted}

\begin{clinicalnote}
The slope of 0.7 means that, on average, systolic BP increases by 0.7~mmHg for each additional year of age. This does \emph{not} mean aging \emph{causes} higher BP --- it means the two are linearly associated in this sample.
\end{clinicalnote}

% ============================================================
\section{Multiple linear regression}

\textbf{Question:} Predict systolic BP from age, BMI, and smoking status simultaneously.

\begin{minted}{python}
import statsmodels.api as sm

# Expanded simulated data
np.random.seed(42)
n = 300
age = np.random.uniform(25, 80, n)
bmi = np.random.normal(27, 5, n)
smoking = np.random.binomial(1, 0.25, n)  # 25% smokers
# True model: SBP = 70 + 0.6*age + 0.8*BMI + 5*smoking + noise
sbp = 70 + 0.6 * age + 0.8 * bmi + 5 * smoking + np.random.normal(0, 10, n)

df = pd.DataFrame({
    "age": age, "bmi": bmi, "smoking": smoking, "systolic_bp": sbp
})

# Fit multiple regression with statsmodels
X = sm.add_constant(df[["age", "bmi", "smoking"]])
model = sm.OLS(df["systolic_bp"], X).fit()
print(model.summary())
\end{minted}

\subsection{Interpreting coefficients in clinical context}

\begin{minted}{python}
# Extract and display coefficients
coef_df = pd.DataFrame({
    "Coefficient": model.params,
    "95% CI lower": model.conf_int()[0],
    "95% CI upper": model.conf_int()[1],
    "p-value": model.pvalues
}).round(4)
print(coef_df)
\end{minted}

Each coefficient answers: ``holding all other variables constant, what is the expected change in SBP for a one-unit increase in this predictor?''

\begin{itemize}
  \item \textbf{Age coefficient $\approx 0.6$}: each additional year of age is associated with 0.6~mmHg higher SBP, \emph{after adjusting for BMI and smoking}.
  \item \textbf{BMI coefficient $\approx 0.8$}: each additional BMI unit is associated with 0.8~mmHg higher SBP, \emph{after adjusting for age and smoking}.
  \item \textbf{Smoking coefficient $\approx 5$}: smokers have, on average, 5~mmHg higher SBP than non-smokers, \emph{after adjusting for age and BMI}.
\end{itemize}

\begin{pythontip}
The phrase ``after adjusting for'' (or ``controlling for'') is the key difference between simple and multiple regression. It is why epidemiologists use multiple regression --- to separate the effect of one variable from potential confounders.
\end{pythontip}

% ============================================================
\section{Confounders}

A confounder is a variable that is associated with both the exposure and the outcome, distorting the apparent relationship between them.

\begin{minted}{python}
# Example: coffee drinking and heart disease
# Both are associated with smoking (the confounder)
np.random.seed(42)
n = 500
smoking = np.random.binomial(1, 0.3, n)
coffee = 1.5 + 2 * smoking + np.random.normal(0, 1.5, n)  # smokers drink more coffee
heart_disease_risk = 0.5 * smoking + np.random.normal(0, 0.3, n)  # smoking causes HD

# Unadjusted: coffee appears to predict heart disease
r_unadjusted, p_unadjusted = stats.pearsonr(coffee, heart_disease_risk)
print(f"Unadjusted correlation (coffee vs HD risk): r = {r_unadjusted:.3f}, p = {p_unadjusted:.4f}")

# Adjusted regression: coffee effect disappears
df_conf = pd.DataFrame({"coffee": coffee, "smoking": smoking, "hd_risk": heart_disease_risk})
X = sm.add_constant(df_conf[["coffee", "smoking"]])
model_adj = sm.OLS(df_conf["hd_risk"], X).fit()
print(f"\nAdjusted model:")
print(f"  Coffee coefficient: {model_adj.params['coffee']:.4f} (p = {model_adj.pvalues['coffee']:.4f})")
print(f"  Smoking coefficient: {model_adj.params['smoking']:.4f} (p = {model_adj.pvalues['smoking']:.4f})")
\end{minted}

\begin{warningbox}
Without adjusting for smoking, you would wrongly conclude that coffee increases heart disease risk. This is one of the most common errors in observational studies. Always think about what confounders might be lurking in your data.
\end{warningbox}

% ============================================================
\section{Checking regression assumptions}

\begin{minted}{python}
# Residual diagnostics for our SBP model
residuals = model.resid
fitted = model.fittedvalues

fig, axes = plt.subplots(1, 3, figsize=(15, 4))

# Residuals vs fitted
axes[0].scatter(fitted, residuals, alpha=0.4, s=15)
axes[0].axhline(0, color="red", linestyle="--")
axes[0].set_xlabel("Fitted values")
axes[0].set_ylabel("Residuals")
axes[0].set_title("Residuals vs Fitted")

# Histogram of residuals
axes[1].hist(residuals, bins=25, edgecolor="black", alpha=0.7)
axes[1].set_xlabel("Residuals")
axes[1].set_title("Distribution of residuals")

# Q-Q plot
stats.probplot(residuals, dist="norm", plot=axes[2])
axes[2].set_title("Q-Q plot")

plt.tight_layout()
plt.show()
\end{minted}

\begin{clinicalnote}
The four key assumptions of linear regression: (1)~linearity, (2)~independence of residuals, (3)~homoscedasticity (constant variance), (4)~normality of residuals. The plots above help you check assumptions 1, 3, and 4. Violation does not always invalidate results, but it should make you cautious.
\end{clinicalnote}

% ============================================================
\section{Logistic regression}

When the outcome is binary (disease: yes/no), linear regression is inappropriate. Logistic regression models the \emph{probability} of the outcome.

\textbf{Question:} Predict whether a patient has diabetes (yes/no) from age, BMI, and glucose level.

\begin{minted}{python}
from sklearn.linear_model import LogisticRegression
from sklearn.model_selection import train_test_split
from sklearn.metrics import classification_report, confusion_matrix

# Load the Pima Indians Diabetes dataset
url = ("https://raw.githubusercontent.com/jbrownlee/Datasets/"
       "master/pima-indians-diabetes.data.csv")
columns = ["pregnancies", "glucose", "blood_pressure", "skin_thickness",
           "insulin", "bmi", "diabetes_pedigree", "age", "outcome"]
pima = pd.read_csv(url, names=columns)
print(f"Shape: {pima.shape}")
print(f"Diabetes prevalence: {pima['outcome'].mean():.1%}")
pima.head()
\end{minted}

\begin{minted}{python}
# Fit logistic regression with statsmodels (for detailed output)
X = sm.add_constant(pima[["age", "bmi", "glucose"]])
y = pima["outcome"]

logit_model = sm.Logit(y, X).fit()
print(logit_model.summary())
\end{minted}

% ============================================================
\section{Odds ratios from logistic regression}

The coefficients of logistic regression are log-odds. Exponentiate them to get \textbf{odds ratios}:

\begin{minted}{python}
# Odds ratios with 95% confidence intervals
odds_ratios = np.exp(logit_model.params)
ci = np.exp(logit_model.conf_int())

or_df = pd.DataFrame({
    "Odds Ratio": odds_ratios,
    "95% CI lower": ci[0],
    "95% CI upper": ci[1],
    "p-value": logit_model.pvalues
}).round(4)
print(or_df)
\end{minted}

\begin{clinicalnote}
An odds ratio of 1.04 for glucose means: for each additional 1~mg/dL of glucose, the odds of diabetes increase by 4\%, holding age and BMI constant. An OR $> 1$ means higher risk; OR $< 1$ means lower risk; OR $= 1$ means no association.
\end{clinicalnote}

% ============================================================
\section{Model evaluation: confusion matrix and beyond}

\begin{minted}{python}
# Fit with scikit-learn for prediction
features = ["age", "bmi", "glucose"]
X = pima[features]
y = pima["outcome"]

X_train, X_test, y_train, y_test = train_test_split(
    X, y, test_size=0.3, random_state=42, stratify=y
)

clf = LogisticRegression(max_iter=1000, random_state=42)
clf.fit(X_train, y_train)
y_pred = clf.predict(X_test)
\end{minted}

\begin{minted}{python}
# Confusion matrix
cm = confusion_matrix(y_test, y_pred)
print("Confusion Matrix:")
print(cm)
print()

# Detailed metrics
print(classification_report(y_test, y_pred,
                            target_names=["No Diabetes", "Diabetes"]))
\end{minted}

\subsection{Sensitivity, specificity, PPV, NPV}

\begin{minted}{python}
tn, fp, fn, tp = cm.ravel()
sensitivity = tp / (tp + fn)    # true positive rate (recall)
specificity = tn / (tn + fp)    # true negative rate
ppv = tp / (tp + fp)            # positive predictive value (precision)
npv = tn / (tn + fn)            # negative predictive value

print(f"Sensitivity (recall): {sensitivity:.3f}")
print(f"  -> Of those WITH diabetes, {sensitivity:.1%} were correctly identified")
print(f"Specificity:          {specificity:.3f}")
print(f"  -> Of those WITHOUT diabetes, {specificity:.1%} were correctly identified")
print(f"PPV (precision):      {ppv:.3f}")
print(f"  -> Of those PREDICTED positive, {ppv:.1%} actually have diabetes")
print(f"NPV:                  {npv:.3f}")
print(f"  -> Of those PREDICTED negative, {npv:.1%} are truly disease-free")
\end{minted}

\begin{minted}{python}
# Visualize the confusion matrix
fig, ax = plt.subplots(figsize=(6, 5))
sns.heatmap(cm, annot=True, fmt="d", cmap="Blues",
            xticklabels=["No Diabetes", "Diabetes"],
            yticklabels=["No Diabetes", "Diabetes"], ax=ax)
ax.set_xlabel("Predicted")
ax.set_ylabel("Actual")
ax.set_title("Confusion Matrix")
plt.tight_layout()
plt.show()
\end{minted}

\begin{warningbox}
In clinical screening, missing a disease (false negative) is often more dangerous than a false alarm (false positive). A screening test for cancer should have \textbf{high sensitivity} even if specificity is moderate. A confirmatory test should have \textbf{high specificity}.
\end{warningbox}

% ============================================================
\section{Introduction to survival analysis}

Sometimes the outcome is not ``yes/no'' but ``time until event'' --- time to death, time to relapse, time to hospital readmission. This is \textbf{survival analysis}.

\subsection{Why not just use logistic regression?}

Two reasons:
\begin{enumerate}
  \item \textbf{Censoring.} Some patients leave the study before the event occurs. We know they survived \emph{at least} that long, but not their true event time.
  \item \textbf{Time matters.} A drug that delays relapse by 3 years is better than one that delays it by 3 months, even if both eventually fail.
\end{enumerate}

\subsection{Kaplan-Meier curves}

\begin{minted}{python}
# Install lifelines if needed: pip install lifelines
from lifelines import KaplanMeierFitter, CoxPHFitter

# Simulated survival data
np.random.seed(42)
n = 200
treatment = np.random.binomial(1, 0.5, n)
# Treatment group survives longer
time = np.where(treatment == 1,
                np.random.exponential(24, n),   # treatment: mean 24 months
                np.random.exponential(16, n))   # control: mean 16 months
event = np.random.binomial(1, 0.75, n)  # 25% censored

surv_df = pd.DataFrame({
    "time": time, "event": event, "treatment": treatment
})
\end{minted}

\begin{minted}{python}
# Kaplan-Meier curves
fig, ax = plt.subplots(figsize=(8, 5))

for label, group_df in surv_df.groupby("treatment"):
    kmf = KaplanMeierFitter()
    kmf.fit(group_df["time"], event_observed=group_df["event"],
            label="Treatment" if label == 1 else "Control")
    kmf.plot_survival_function(ax=ax)

ax.set_xlabel("Time (months)")
ax.set_ylabel("Survival probability")
ax.set_title("Kaplan-Meier survival curves")
plt.tight_layout()
plt.show()
\end{minted}

\subsection{Cox proportional hazards model}

The Cox model is the ``regression'' of survival analysis. It estimates how covariates affect the hazard (risk of event at each time point):

\begin{minted}{python}
# Cox proportional hazards
cph = CoxPHFitter()
cph.fit(surv_df, duration_col="time", event_col="event",
        formula="treatment")
cph.print_summary()
\end{minted}

\begin{minted}{python}
# Hazard ratio
hr = np.exp(cph.params_["treatment"])
print(f"Hazard ratio for treatment: {hr:.3f}")
if hr < 1:
    print(f"Treatment reduces hazard by {(1-hr)*100:.1f}%")
else:
    print(f"Treatment increases hazard by {(hr-1)*100:.1f}%")
\end{minted}

\begin{clinicalnote}
A hazard ratio of 0.65 means the treatment group has 35\% lower risk of the event at any given time point, compared to the control group. Hazard ratios are the standard way to report results in oncology trials, cardiovascular outcome studies, and any study with a time-to-event endpoint.
\end{clinicalnote}

% ============================================================
\section{Practical: diabetes risk prediction}

\begin{exercise}
Using the Pima Indians Diabetes dataset (\url{https://raw.githubusercontent.com/jbrownlee/Datasets/master/pima-indians-diabetes.data.csv}):
\begin{enumerate}
  \item \textbf{Data preparation.} Load the dataset. Replace zero values in \texttt{glucose}, \texttt{blood\_pressure}, \texttt{bmi}, \texttt{skin\_thickness}, and \texttt{insulin} with \texttt{NaN} (these are biologically impossible zeros). Drop rows with missing values.
  \item \textbf{Simple logistic regression.} Fit a model predicting diabetes from glucose alone. What is the odds ratio? Interpret it.
  \item \textbf{Multiple logistic regression.} Fit a model using age, BMI, glucose, blood pressure, and diabetes pedigree function. Which variables are statistically significant? Report odds ratios with 95\% CIs.
  \item \textbf{Model comparison.} Split the data 70/30. Fit both models on the training set. Compare sensitivity, specificity, and accuracy on the test set. Which model is better for \emph{screening}?
  \item \textbf{Confounding.} Fit a model with glucose only, then add age. Does the glucose coefficient change? What does this tell you about confounding?
  \item \textbf{Clinical thresholds.} Instead of the default 0.5 threshold, try 0.3 and 0.7. How do sensitivity and specificity change? When might you prefer a lower threshold?
  \item \textbf{Bonus: survival analysis.} Using the simulated survival data above, add ``age'' and ``smoking'' as covariates to the Cox model. Which has the largest hazard ratio?
\end{enumerate}
\end{exercise}

% ============================================================
\section{Chapter summary}

\begin{itemize}
  \item Simple linear regression models the relationship between one predictor and a continuous outcome.
  \item Multiple regression adjusts for confounders --- the phrase ``after controlling for'' reflects this.
  \item Coefficients represent the expected change in the outcome for a one-unit change in the predictor, holding other variables constant.
  \item Logistic regression predicts binary outcomes (disease yes/no) and produces odds ratios.
  \item Model evaluation requires the confusion matrix: sensitivity, specificity, PPV, and NPV.
  \item High sensitivity is critical for screening; high specificity is critical for confirmation.
  \item Survival analysis (Kaplan-Meier, Cox model) handles time-to-event data with censoring.
  \item Always check regression assumptions and report confidence intervals, not just p-values.
\end{itemize}

\chapter{Machine learning for clinical prediction}
\label{ch:ml-clinical}

\begin{quote}
\emph{``A model that works perfectly on your training data and fails on the next patient is not a model --- it is a memory.''}
\end{quote}

% ============================================================
\section{When to use ML vs traditional statistics}

Machine learning and traditional statistics answer different questions:

\begin{itemize}
  \item \textbf{Traditional statistics} (regression, t-tests): ``Which risk factors are associated with heart disease, and by how much?'' The goal is \emph{inference} --- understanding relationships.
  \item \textbf{Machine learning}: ``Given this patient's data, what is their probability of heart disease?'' The goal is \emph{prediction} --- accuracy on new patients.
\end{itemize}

\begin{clinicalnote}
Use logistic regression when you need to explain \emph{why} (odds ratios, confidence intervals, p-values for each risk factor). Use ML when you need the best possible \emph{prediction} and interpretability is secondary. In practice, many clinical studies use both: regression to understand the biology, ML to build the screening tool.
\end{clinicalnote}

\subsection{When NOT to use ML}

\begin{itemize}
  \item Small datasets ($< 200$ patients): ML overfits easily. Stick with logistic regression.
  \item When you need regulatory approval: the FDA requires interpretable models for most clinical decision support tools.
  \item When a simple rule works: if ``glucose $> 126$'' classifies 90\% of diabetics correctly, you do not need a random forest.
\end{itemize}

\begin{minted}{python}
import numpy as np
import pandas as pd
import matplotlib.pyplot as plt
import seaborn as sns
from sklearn.model_selection import train_test_split, cross_val_score
from sklearn.preprocessing import StandardScaler
from sklearn.linear_model import LogisticRegression
from sklearn.tree import DecisionTreeClassifier, plot_tree
from sklearn.ensemble import RandomForestClassifier
from sklearn.metrics import (classification_report, confusion_matrix,
                             roc_curve, roc_auc_score, RocCurveDisplay)
\end{minted}

% ============================================================
\section{Loading clinical data}

We use the UCI Heart Disease dataset --- a classic benchmark in clinical ML.

\begin{minted}{python}
# UCI Heart Disease dataset
url = ("https://raw.githubusercontent.com/raphaelfontenelle/"
       "heart-disease-uci/main/heart.csv")
heart = pd.read_csv(url)
print(f"Shape: {heart.shape}")
print(f"Heart disease prevalence: {heart['target'].mean():.1%}")
heart.head()
\end{minted}

\begin{minted}{python}
# Column descriptions
column_info = {
    "age": "Age in years",
    "sex": "1 = male, 0 = female",
    "cp": "Chest pain type (0-3)",
    "trestbps": "Resting blood pressure (mmHg)",
    "chol": "Serum cholesterol (mg/dL)",
    "fbs": "Fasting blood sugar > 120 mg/dL (1 = true)",
    "restecg": "Resting ECG results (0-2)",
    "thalach": "Maximum heart rate achieved",
    "exang": "Exercise-induced angina (1 = yes)",
    "oldpeak": "ST depression induced by exercise",
    "slope": "Slope of peak exercise ST segment",
    "ca": "Number of major vessels colored by fluoroscopy (0-3)",
    "thal": "Thalassemia (1 = normal, 2 = fixed defect, 3 = reversible)",
    "target": "Heart disease (1 = yes, 0 = no)"
}
for col, desc in column_info.items():
    print(f"  {col:12s} : {desc}")
\end{minted}

% ============================================================
\section{Train/test split and cross-validation}

\subsection{Why split the data?}

If you train a model on all your data and then evaluate it on the same data, the accuracy is misleadingly high. The model has ``seen the answers.'' To estimate real-world performance, you must evaluate on data the model has \emph{never seen}.

\begin{minted}{python}
# Define features and target
feature_cols = ["age", "sex", "cp", "trestbps", "chol", "fbs",
                "restecg", "thalach", "exang", "oldpeak", "slope", "ca", "thal"]
X = heart[feature_cols]
y = heart["target"]

# 70/30 train/test split (stratified to preserve class balance)
X_train, X_test, y_train, y_test = train_test_split(
    X, y, test_size=0.3, random_state=42, stratify=y
)
print(f"Training set: {X_train.shape[0]} patients")
print(f"Test set:     {X_test.shape[0]} patients")
print(f"Training prevalence: {y_train.mean():.1%}")
print(f"Test prevalence:     {y_test.mean():.1%}")
\end{minted}

\subsection{Cross-validation}

A single train/test split can be lucky or unlucky. Cross-validation (CV) gives a more robust estimate by rotating the test set:

\begin{minted}{python}
# 5-fold cross-validation with logistic regression
lr = LogisticRegression(max_iter=1000, random_state=42)
cv_scores = cross_val_score(lr, X, y, cv=5, scoring="accuracy")
print(f"CV accuracy: {cv_scores.mean():.3f} +/- {cv_scores.std():.3f}")
print(f"Per-fold: {cv_scores.round(3)}")
\end{minted}

\begin{warningbox}
Never use test set results to choose your model or tune hyperparameters. The test set must only be used \textbf{once}, at the very end. If you peek at test results during model selection, you are effectively training on the test set.
\end{warningbox}

% ============================================================
\section{Decision trees}

A decision tree splits patients into groups using yes/no questions, forming a flowchart that clinicians can follow.

\begin{minted}{python}
# Fit a decision tree
dt = DecisionTreeClassifier(max_depth=4, random_state=42)
dt.fit(X_train, y_train)

# Evaluate
y_pred_dt = dt.predict(X_test)
print("Decision Tree Results:")
print(classification_report(y_test, y_pred_dt,
                            target_names=["No Disease", "Disease"]))
\end{minted}

\begin{minted}{python}
# Visualize the tree
fig, ax = plt.subplots(figsize=(20, 10))
plot_tree(dt, feature_names=feature_cols,
          class_names=["No Disease", "Disease"],
          filled=True, rounded=True, fontsize=9, ax=ax)
plt.title("Decision tree for heart disease prediction")
plt.tight_layout()
plt.show()
\end{minted}

\begin{clinicalnote}
Decision trees are popular in emergency medicine because they produce \emph{clinical decision rules} --- simple flowcharts that a physician can follow at the bedside without a computer. The HEART score, Wells criteria, and Ottawa ankle rules are all conceptually similar to decision trees.
\end{clinicalnote}

\subsection{The danger of deep trees}

\begin{minted}{python}
# Overfit tree (no depth limit)
dt_overfit = DecisionTreeClassifier(random_state=42)  # no max_depth
dt_overfit.fit(X_train, y_train)

print(f"Training accuracy (overfit): {dt_overfit.score(X_train, y_train):.3f}")
print(f"Test accuracy (overfit):     {dt_overfit.score(X_test, y_test):.3f}")
print(f"Training accuracy (depth=4): {dt.score(X_train, y_train):.3f}")
print(f"Test accuracy (depth=4):     {dt.score(X_test, y_test):.3f}")
\end{minted}

\begin{pythontip}
When training accuracy is much higher than test accuracy, your model is \textbf{overfitting} --- it has memorized the training data instead of learning general patterns. Limit tree depth with \texttt{max\_depth} or require a minimum number of samples per leaf with \texttt{min\_samples\_leaf}.
\end{pythontip}

% ============================================================
\section{Random forests}

A random forest combines hundreds of decision trees, each trained on a random subset of the data and features. The final prediction is the majority vote.

\begin{minted}{python}
# Fit a random forest
rf = RandomForestClassifier(n_estimators=200, max_depth=6,
                            random_state=42, n_jobs=-1)
rf.fit(X_train, y_train)

y_pred_rf = rf.predict(X_test)
print("Random Forest Results:")
print(classification_report(y_test, y_pred_rf,
                            target_names=["No Disease", "Disease"]))
\end{minted}

\begin{minted}{python}
# Cross-validation comparison
models = {
    "Logistic Regression": LogisticRegression(max_iter=1000, random_state=42),
    "Decision Tree (depth=4)": DecisionTreeClassifier(max_depth=4, random_state=42),
    "Random Forest": RandomForestClassifier(n_estimators=200, max_depth=6,
                                            random_state=42, n_jobs=-1)
}

print("5-fold CV accuracy:")
for name, model in models.items():
    scores = cross_val_score(model, X, y, cv=5, scoring="accuracy")
    print(f"  {name:30s}: {scores.mean():.3f} +/- {scores.std():.3f}")
\end{minted}

% ============================================================
\section{ROC curves and AUC}

The ROC (Receiver Operating Characteristic) curve plots sensitivity vs (1 -- specificity) across all possible classification thresholds. The AUC (Area Under the Curve) summarizes discriminative ability in a single number.

\begin{itemize}
  \item AUC = 0.5: no better than flipping a coin.
  \item AUC = 0.7--0.8: acceptable discrimination.
  \item AUC = 0.8--0.9: excellent discrimination.
  \item AUC $> 0.9$: outstanding (rare in clinical practice).
\end{itemize}

\begin{minted}{python}
# Get predicted probabilities
lr.fit(X_train, y_train)
y_prob_lr = lr.predict_proba(X_test)[:, 1]
y_prob_rf = rf.predict_proba(X_test)[:, 1]

# Compute ROC curves
fpr_lr, tpr_lr, _ = roc_curve(y_test, y_prob_lr)
fpr_rf, tpr_rf, _ = roc_curve(y_test, y_prob_rf)
auc_lr = roc_auc_score(y_test, y_prob_lr)
auc_rf = roc_auc_score(y_test, y_prob_rf)

# Plot
fig, ax = plt.subplots(figsize=(7, 6))
ax.plot(fpr_lr, tpr_lr, label=f"Logistic Regression (AUC = {auc_lr:.3f})",
        linewidth=2)
ax.plot(fpr_rf, tpr_rf, label=f"Random Forest (AUC = {auc_rf:.3f})",
        linewidth=2)
ax.plot([0, 1], [0, 1], "k--", alpha=0.5, label="Random (AUC = 0.500)")
ax.set_xlabel("False Positive Rate (1 - Specificity)")
ax.set_ylabel("True Positive Rate (Sensitivity)")
ax.set_title("ROC curves: Heart disease prediction")
ax.legend(loc="lower right")
plt.tight_layout()
plt.show()
\end{minted}

\begin{clinicalnote}
In clinical practice, ROC curves are used to compare diagnostic tests. When evaluating a new biomarker for sepsis, you plot its ROC curve against existing biomarkers (CRP, procalcitonin). The test with the highest AUC has the best overall discriminative ability.
\end{clinicalnote}

% ============================================================
\section{Feature importance}

\textbf{Question:} Which risk factors matter most for predicting heart disease?

\begin{minted}{python}
# Random forest feature importance
importances = rf.feature_importances_
importance_df = pd.DataFrame({
    "Feature": feature_cols,
    "Importance": importances
}).sort_values("Importance", ascending=True)

fig, ax = plt.subplots(figsize=(8, 6))
ax.barh(importance_df["Feature"], importance_df["Importance"], color="steelblue")
ax.set_xlabel("Feature importance (Gini)")
ax.set_title("Which risk factors predict heart disease?")
plt.tight_layout()
plt.show()
\end{minted}

\begin{minted}{python}
# Top 5 features
print("Top 5 most important features:")
for _, row in importance_df.tail(5).iloc[::-1].iterrows():
    print(f"  {row['Feature']:12s}: {row['Importance']:.3f}")
\end{minted}

\begin{warningbox}
Feature importance from random forests tells you which variables are useful for \emph{prediction}, not which variables \emph{cause} the outcome. A variable can be important for prediction because it is a proxy for something else. Do not confuse predictive importance with causal importance.
\end{warningbox}

% ============================================================
\section{Calibration}

A model says ``this patient has a 70\% risk of heart disease.'' Does that mean that among all patients the model gives 70\%, roughly 70\% actually have heart disease? If yes, the model is \textbf{well-calibrated}.

\begin{minted}{python}
from sklearn.calibration import calibration_curve

# Calibration plot for random forest
prob_true, prob_pred = calibration_curve(y_test, y_prob_rf, n_bins=8)

fig, ax = plt.subplots(figsize=(6, 6))
ax.plot(prob_pred, prob_true, "o-", linewidth=2, label="Random Forest")
ax.plot([0, 1], [0, 1], "k--", label="Perfectly calibrated")
ax.set_xlabel("Mean predicted probability")
ax.set_ylabel("Observed proportion")
ax.set_title("Calibration plot")
ax.legend()
plt.tight_layout()
plt.show()
\end{minted}

\begin{minted}{python}
# Brier score: lower is better (0 = perfect)
from sklearn.metrics import brier_score_loss

brier_lr = brier_score_loss(y_test, y_prob_lr)
brier_rf = brier_score_loss(y_test, y_prob_rf)
print(f"Brier score (Logistic Regression): {brier_lr:.4f}")
print(f"Brier score (Random Forest):       {brier_rf:.4f}")
\end{minted}

\begin{clinicalnote}
Calibration is critical in clinical prediction models. If a diabetes risk calculator tells a patient ``you have a 40\% chance of developing diabetes in 10 years,'' that number must be accurate --- it drives treatment decisions. The Framingham Risk Score and QRISK are regularly re-calibrated for different populations.
\end{clinicalnote}

% ============================================================
\section{Overfitting: the danger of small clinical datasets}

\begin{minted}{python}
# Demonstrate overfitting with varying dataset sizes
train_sizes = [30, 50, 100, 150, 200, len(X_train)]
results = []

for size in train_sizes:
    if size > len(X_train):
        continue
    X_sub = X_train.iloc[:size]
    y_sub = y_train.iloc[:size]

    rf_temp = RandomForestClassifier(n_estimators=100, random_state=42)
    rf_temp.fit(X_sub, y_sub)

    train_acc = rf_temp.score(X_sub, y_sub)
    test_acc = rf_temp.score(X_test, y_test)
    results.append({"n_train": size, "train_acc": train_acc, "test_acc": test_acc})

results_df = pd.DataFrame(results)

fig, ax = plt.subplots(figsize=(8, 5))
ax.plot(results_df["n_train"], results_df["train_acc"], "o-",
        label="Training accuracy", linewidth=2)
ax.plot(results_df["n_train"], results_df["test_acc"], "s-",
        label="Test accuracy", linewidth=2)
ax.set_xlabel("Number of training patients")
ax.set_ylabel("Accuracy")
ax.set_title("Learning curve: more data reduces overfitting")
ax.legend()
ax.set_ylim(0.5, 1.05)
plt.tight_layout()
plt.show()
\end{minted}

\begin{pythontip}
Signs of overfitting: (1)~training accuracy is much higher than test accuracy, (2)~performance varies wildly across cross-validation folds, (3)~adding more features hurts test performance. The cure: more data, simpler models, regularization, or cross-validation for model selection.
\end{pythontip}

% ============================================================
\section{Practical: heart disease prediction pipeline}

\begin{minted}{python}
# Complete pipeline from raw data to evaluated model
from sklearn.pipeline import Pipeline
from sklearn.impute import SimpleImputer

# Full pipeline
pipeline = Pipeline([
    ("imputer", SimpleImputer(strategy="median")),
    ("scaler", StandardScaler()),
    ("classifier", RandomForestClassifier(n_estimators=200, max_depth=6,
                                          random_state=42))
])

# Cross-validate the full pipeline
cv_scores = cross_val_score(pipeline, X, y, cv=5, scoring="roc_auc")
print(f"Pipeline CV AUC: {cv_scores.mean():.3f} +/- {cv_scores.std():.3f}")

# Final fit and evaluation
pipeline.fit(X_train, y_train)
y_prob_final = pipeline.predict_proba(X_test)[:, 1]
y_pred_final = pipeline.predict(X_test)

print(f"\nFinal test AUC: {roc_auc_score(y_test, y_prob_final):.3f}")
print(f"\n{classification_report(y_test, y_pred_final, target_names=['No Disease', 'Disease'])}")
\end{minted}

% ============================================================
\section{Ethics: fairness in clinical prediction models}

\begin{warningbox}
Clinical prediction models can perpetuate and amplify health disparities. A model trained predominantly on one population may perform poorly on another. Before deploying any clinical ML model, ask:
\begin{itemize}
  \item \textbf{Who is in the training data?} If 90\% of patients are white males, the model may not work for women or minority groups.
  \item \textbf{Does performance differ across subgroups?} Check AUC, sensitivity, and specificity separately by sex, race, and age group.
  \item \textbf{What are the consequences of errors?} A false negative for cardiac risk in a young woman is more dangerous than a false positive.
  \item \textbf{Is the model transparent?} Clinicians and patients deserve to know how a prediction was made.
\end{itemize}
\end{warningbox}

\begin{minted}{python}
# Check model performance by sex
for sex_val, sex_label in [(1, "Male"), (0, "Female")]:
    mask = X_test["sex"] == sex_val
    if mask.sum() < 10:
        print(f"{sex_label}: too few samples ({mask.sum()})")
        continue

    auc_sub = roc_auc_score(y_test[mask], y_prob_final[mask])
    sens = (y_pred_final[mask] == 1)[y_test[mask] == 1].mean()
    spec = (y_pred_final[mask] == 0)[y_test[mask] == 0].mean()

    print(f"{sex_label:8s}: AUC = {auc_sub:.3f}, "
          f"Sensitivity = {sens:.3f}, Specificity = {spec:.3f}, "
          f"n = {mask.sum()}")
\end{minted}

\begin{clinicalnote}
In 2019, a widely-used commercial algorithm for allocating healthcare resources was found to be biased against Black patients. The algorithm used healthcare \emph{costs} as a proxy for healthcare \emph{needs}, but because Black patients had lower costs due to systemic barriers to access, the algorithm systematically underestimated their needs. This affected millions of patients. The lesson: always audit your model for disparities.
\end{clinicalnote}

\begin{exercise}
Using the UCI Heart Disease dataset:
\begin{enumerate}
  \item \textbf{Baseline.} Fit a logistic regression model using all features. Report CV accuracy and AUC.
  \item \textbf{Decision tree.} Fit a decision tree with \texttt{max\_depth=3}. Visualize the tree. Can a clinician follow this decision rule at the bedside?
  \item \textbf{Random forest.} Fit a random forest with 200 trees. Compare its AUC to logistic regression and the decision tree. Is the improvement worth the loss of interpretability?
  \item \textbf{Feature selection.} Train a random forest using only the top 5 most important features. How does performance compare to using all 13 features?
  \item \textbf{ROC comparison.} Plot ROC curves for all three models on the same figure. Which model would you choose for a screening tool? For a confirmatory test?
  \item \textbf{Threshold analysis.} For the random forest, plot sensitivity and specificity as a function of the classification threshold (from 0 to 1). At what threshold are they equal?
  \item \textbf{Calibration.} Plot calibration curves for logistic regression and random forest. Which model is better calibrated?
  \item \textbf{Fairness audit.} Compare AUC, sensitivity, and specificity between male and female patients. Are there concerning disparities? What might explain them?
  \item \textbf{Clinical deployment.} Write a paragraph (in a Markdown cell) describing how you would validate this model before using it in a real hospital. Consider: external validation, regulatory requirements, clinician workflow, and patient consent.
\end{enumerate}
\end{exercise}

% ============================================================
\section{Chapter summary}

\begin{itemize}
  \item Use traditional statistics (regression) for inference and ML for prediction.
  \item Always split data into training and test sets. Never evaluate on training data.
  \item Cross-validation gives more robust performance estimates than a single split.
  \item Decision trees are interpretable and clinician-friendly but prone to overfitting.
  \item Random forests improve accuracy by averaging many trees.
  \item ROC curves and AUC compare the discriminative ability of different models.
  \item Feature importance reveals which risk factors drive predictions (but not causation).
  \item Calibration checks whether predicted probabilities match observed outcomes.
  \item Overfitting is the central danger of ML on small clinical datasets --- always validate.
  \item Fairness audits are essential: check model performance across demographic subgroups before deployment.
\end{itemize}

\chapter{Geospatial and temporal health data}
\label{ch:geospatial-temporal}

\begin{quote}
\emph{``An epidemic is a story told in time and space. If you can plot both, you can see where it came from and where it is going.''}
\end{quote}

% ============================================================
\section{Time series in health}

A \emph{time series} is a sequence of measurements recorded at regular intervals: daily case counts, weekly death tolls, monthly vaccination doses. Health time series have three components:

\begin{itemize}
  \item \textbf{Trend.} The long-term direction --- is incidence rising or falling?
  \item \textbf{Seasonality.} Regular cycles --- malaria peaks during rainy seasons, influenza peaks in winter.
  \item \textbf{Noise.} Random day-to-day fluctuation that obscures the signal.
\end{itemize}

\begin{clinicalnote}
Epidemic curves (``epi curves'') are the most important time series in public health. During an outbreak, the shape of the epi curve tells you whether the source is a point exposure (sharp peak), a continuous source (plateau), or person-to-person transmission (successive waves).
\end{clinicalnote}

% ============================================================
\section{Plotting time series with pandas}

\subsection{Date parsing and indexing}

Pandas handles dates natively. The key is to parse date columns and set them as the index:

\begin{minted}{python}
import pandas as pd
import matplotlib.pyplot as plt

# Example: daily malaria cases at a district hospital
dates = pd.date_range("2023-01-01", periods=365, freq="D")
import numpy as np
np.random.seed(42)
# Simulate seasonal pattern: peak during rainy season (June-September)
day_of_year = dates.dayofyear
seasonal = 50 + 40 * np.sin(2 * np.pi * (day_of_year - 90) / 365)
cases = np.maximum(0, seasonal + np.random.normal(0, 12, 365)).astype(int)

df = pd.DataFrame({"date": dates, "cases": cases})
df["date"] = pd.to_datetime(df["date"])
df = df.set_index("date")
df.head()
\end{minted}

\subsection{Rolling averages}

Day-to-day fluctuation makes trends hard to see. A \emph{rolling average} smooths the curve:

\begin{minted}{python}
df["cases_7d"] = df["cases"].rolling(window=7).mean()
df["cases_14d"] = df["cases"].rolling(window=14).mean()

fig, ax = plt.subplots(figsize=(12, 5))
ax.bar(df.index, df["cases"], alpha=0.3, label="Daily cases", color="steelblue")
ax.plot(df.index, df["cases_7d"], color="red", linewidth=2, label="7-day average")
ax.plot(df.index, df["cases_14d"], color="darkgreen", linewidth=2, label="14-day average")
ax.set_xlabel("Date")
ax.set_ylabel("Malaria cases")
ax.set_title("Daily malaria cases with rolling averages")
ax.legend()
plt.tight_layout()
plt.show()
\end{minted}

\begin{pythontip}
The \texttt{window} parameter sets the number of days to average over. A 7-day window removes day-of-week effects; a 14-day window smooths more aggressively. Choose based on how noisy your data is.
\end{pythontip}

\subsection{Trend decomposition}

The \texttt{statsmodels} library can decompose a time series into trend, seasonal, and residual components:

\begin{minted}{python}
from statsmodels.tsa.seasonal import seasonal_decompose

result = seasonal_decompose(df["cases"], model="additive", period=30)

fig, axes = plt.subplots(4, 1, figsize=(12, 10), sharex=True)
result.observed.plot(ax=axes[0], title="Observed")
result.trend.plot(ax=axes[1], title="Trend")
result.seasonal.plot(ax=axes[2], title="Seasonal")
result.resid.plot(ax=axes[3], title="Residual")
plt.tight_layout()
plt.show()
\end{minted}

\begin{warningbox}
The \texttt{period} parameter must match the expected cycle length. For daily data with monthly seasonality, use \texttt{period=30}. For weekly data with annual seasonality, use \texttt{period=52}. A wrong period will produce misleading decompositions.
\end{warningbox}

% ============================================================
\section{COVID-19 case curves}

\subsection{Loading JHU data}

The Johns Hopkins University CSSE COVID-19 dataset is one of the most cited public health datasets in history:

\begin{minted}{python}
# JHU CSSE COVID-19 confirmed cases (global, time series)
url = ("https://raw.githubusercontent.com/CSSEGISandData/COVID-19/"
       "master/csse_covid_19_data/csse_covid_19_time_series/"
       "time_series_covid19_confirmed_global.csv")
confirmed = pd.read_csv(url)
print(f"Shape: {confirmed.shape}")
confirmed.head()
\end{minted}

\subsection{Reshaping wide to long format}

The JHU data is in wide format (one column per date). We need to reshape it:

\begin{minted}{python}
# Focus on 5 African countries
countries = ["South Africa", "Nigeria", "Kenya", "Egypt", "Morocco"]
africa = confirmed[confirmed["Country/Region"].isin(countries)]

# Group by country (some countries have province-level rows)
africa = africa.groupby("Country/Region").sum(numeric_only=True)
africa = africa.drop(columns=["Lat", "Long"], errors="ignore")

# Reshape: wide -> long
africa_long = africa.T
africa_long.index = pd.to_datetime(africa_long.index)
africa_long.index.name = "date"
africa_long.head()
\end{minted}

\subsection{Computing daily cases and 7-day averages}

\begin{minted}{python}
# Cumulative -> daily new cases
daily = africa_long.diff().clip(lower=0)  # clip removes negative corrections

# 7-day rolling average
daily_7d = daily.rolling(7).mean()

# Plot
fig, ax = plt.subplots(figsize=(14, 6))
for country in countries:
    ax.plot(daily_7d.index, daily_7d[country], linewidth=2, label=country)
ax.set_xlabel("Date")
ax.set_ylabel("Daily new cases (7-day average)")
ax.set_title("COVID-19 daily cases in 5 African countries")
ax.legend()
plt.tight_layout()
plt.show()
\end{minted}

\begin{clinicalnote}
The \texttt{.diff()} method converts cumulative counts to daily increments. Negative values occasionally appear because of data corrections (e.g., a country revises its total downward). Clipping at zero is the standard approach, though it slightly inflates subsequent days.
\end{clinicalnote}

% ============================================================
\section{Introduction to geospatial data}

Geospatial data links measurements to locations on Earth. Two common formats:

\begin{itemize}
  \item \textbf{Shapefiles} (\texttt{.shp}) --- the traditional GIS format. A shapefile is actually a set of files (\texttt{.shp}, \texttt{.shx}, \texttt{.dbf}, \texttt{.prj}) that together define polygon geometries and attributes.
  \item \textbf{GeoJSON} (\texttt{.geojson}) --- a single JSON file. Easier to share, read, and use on the web.
\end{itemize}

Both formats store \emph{geometries} (country borders, district boundaries, hospital point locations) and \emph{attributes} (population, case count, mortality rate) in the same structure.

% ============================================================
\section{geopandas for disease mapping}

\texttt{geopandas} extends pandas with spatial capabilities. A \texttt{GeoDataFrame} is a DataFrame with a special \texttt{geometry} column that holds shapes.

\subsection{Installation}

\begin{minted}{python}
# In Google Colab or local environment
!pip install geopandas mapclassify
\end{minted}

\subsection{Loading a world map}

\begin{minted}{python}
import geopandas as gpd

# Natural Earth dataset (built into geopandas)
world = gpd.read_file(gpd.datasets.get_path("naturalearth_lowres"))
print(f"Columns: {world.columns.tolist()}")
print(f"CRS: {world.crs}")  # Coordinate Reference System
world.head()
\end{minted}

\subsection{Choropleth maps}

A \emph{choropleth} map colors regions according to a variable. This is the standard tool for mapping disease burden by country or district:

\begin{minted}{python}
# Filter to Africa
africa_map = world[world["continent"] == "Africa"].copy()

# Plot population estimate
fig, ax = plt.subplots(figsize=(10, 10))
africa_map.plot(column="pop_est", cmap="YlOrRd", legend=True,
                legend_kwds={"label": "Population", "shrink": 0.6},
                edgecolor="black", linewidth=0.5, ax=ax)
ax.set_title("Population estimates in African countries")
ax.set_axis_off()
plt.tight_layout()
plt.show()
\end{minted}

\begin{pythontip}
Good colormaps for health data: \texttt{"YlOrRd"} (sequential, low-to-high severity), \texttt{"RdYlGn\_r"} (diverging, red = bad, green = good), \texttt{"Blues"} (sequential, neutral). Avoid rainbow colormaps (\texttt{"jet"}) --- they mislead the eye.
\end{pythontip}

% ============================================================
\section{Mapping malaria prevalence by African country}

Let us combine real health data with geospatial plotting. We will use the WHO Global Health Observatory malaria incidence data:

\begin{minted}{python}
# WHO GHO: Estimated malaria incidence (per 1000 population at risk)
url = "https://apps.who.int/gho/athena/api/GHO/MALARIA_INCIDENCE?format=csv"
malaria = pd.read_csv(url)
print(malaria.columns.tolist())
malaria.head()
\end{minted}

\begin{minted}{python}
# Filter to most recent year and country-level data
latest_year = malaria["YEAR (DISPLAY)"].max()
malaria_latest = malaria[malaria["YEAR (DISPLAY)"] == latest_year].copy()

# Keep relevant columns
malaria_latest = malaria_latest[["COUNTRY (DISPLAY)", "Numeric"]].copy()
malaria_latest.columns = ["country", "incidence_per_1000"]
malaria_latest.head()
\end{minted}

\begin{minted}{python}
# Merge with Africa map
# Match country names (some need manual fixing)
name_map = {
    "United Republic of Tanzania": "Tanzania",
    "Democratic Republic of the Congo": "Dem. Rep. Congo",
    "Central African Republic": "Central African Rep.",
    "South Sudan": "S. Sudan",
    "Côte d'Ivoire": "Côte d'Ivoire",
    "Equatorial Guinea": "Eq. Guinea",
}
malaria_latest["country"] = malaria_latest["country"].replace(name_map)

africa_malaria = africa_map.merge(malaria_latest, left_on="name",
                                   right_on="country", how="left")
\end{minted}

\begin{minted}{python}
fig, ax = plt.subplots(figsize=(10, 10))
africa_malaria.plot(column="incidence_per_1000", cmap="YlOrRd",
                     legend=True, missing_kwds={"color": "lightgrey"},
                     legend_kwds={"label": "Incidence per 1,000",
                                  "shrink": 0.6},
                     edgecolor="black", linewidth=0.5, ax=ax)
ax.set_title(f"Malaria incidence in Africa ({latest_year})")
ax.set_axis_off()
plt.tight_layout()
plt.show()
\end{minted}

\begin{clinicalnote}
The malaria burden in sub-Saharan Africa accounts for approximately 95\% of global malaria deaths. Choropleth maps immediately reveal the geographic concentration and help target intervention programs such as insecticide-treated nets (ITNs) and indoor residual spraying (IRS).
\end{clinicalnote}

% ============================================================
\section{Combining temporal and spatial analysis}

\subsection{Small multiples: one map per time period}

A practical alternative to animation is the \emph{small multiples} approach --- one map per time period side by side:

\begin{minted}{python}
# Example: COVID-19 cumulative cases per country at 3 time points
dates_to_plot = ["2020-06-01", "2021-01-01", "2021-06-01"]

fig, axes = plt.subplots(1, 3, figsize=(18, 7))
for ax, date_str in zip(axes, dates_to_plot):
    date_col = pd.to_datetime(date_str).strftime("%-m/%-d/%y")
    if date_col in africa.columns:
        data_col = africa[date_col]
    else:
        # Find nearest date
        all_dates = pd.to_datetime(africa.columns, errors="coerce")
        nearest = all_dates[all_dates <= pd.to_datetime(date_str)][-1]
        data_col = africa[nearest.strftime("%-m/%-d/%y")]

    temp = africa_map.copy()
    temp = temp.merge(data_col.reset_index(), left_on="name",
                       right_on="Country/Region", how="left")
    temp.plot(column=data_col.name, cmap="Reds", legend=True,
              edgecolor="black", linewidth=0.5, ax=ax,
              missing_kwds={"color": "lightgrey"})
    ax.set_title(date_str)
    ax.set_axis_off()

plt.suptitle("COVID-19 cumulative cases in Africa", fontsize=14)
plt.tight_layout()
plt.show()
\end{minted}

\subsection{Animated maps (optional advanced)}

For presentations, animated GIFs can show epidemic spread over time. The idea: create one frame per time step, then combine with \texttt{imageio} or \texttt{matplotlib.animation}:

\begin{minted}{python}
import matplotlib.animation as animation

fig, ax = plt.subplots(figsize=(10, 10))

def update(frame_date):
    ax.clear()
    # Merge case data for this date with the map
    temp = africa_map.copy()
    # ... merge logic similar to above ...
    temp.plot(column="cases", cmap="Reds", ax=ax, edgecolor="black",
              linewidth=0.5, missing_kwds={"color": "lightgrey"},
              vmin=0, vmax=max_cases)
    ax.set_title(f"COVID-19 cases — {frame_date}")
    ax.set_axis_off()

# Create animation (one frame per month)
# ani = animation.FuncAnimation(fig, update, frames=monthly_dates,
#                                interval=500)
# ani.save("covid_africa.gif", writer="pillow")
\end{minted}

\begin{warningbox}
Animated maps are visually impressive but can be hard to interpret scientifically. For publication, small multiples or static snapshots with a time slider (in a dashboard) are preferred. Use animations for presentations and communication with the public.
\end{warningbox}

% ============================================================
\section{Practical: COVID-19 dashboard for 5 African countries}

\begin{exercise}
Build a multi-panel COVID-19 dashboard for South Africa, Nigeria, Kenya, Egypt, and Morocco. Your dashboard should have four panels:

\textbf{Panel 1 --- Confirmed cases over time.}
\begin{enumerate}
  \item Load the JHU CSSE time series data for confirmed cases:\\
  \url{https://raw.githubusercontent.com/CSSEGISandData/COVID-19/master/csse_covid_19_data/csse_covid_19_time_series/time_series_covid19_confirmed_global.csv}
  \item Filter to the 5 countries. Compute daily new cases and 7-day rolling averages.
  \item Plot all 5 countries on one axis with a legend.
\end{enumerate}

\textbf{Panel 2 --- Deaths over time.}
\begin{enumerate}
  \item Load the JHU deaths time series:\\
  \url{https://raw.githubusercontent.com/CSSEGISandData/COVID-19/master/csse_covid_19_data/csse_covid_19_time_series/time_series_covid19_deaths_global.csv}
  \item Compute daily deaths (7-day average) and plot.
\end{enumerate}

\textbf{Panel 3 --- Vaccination progress.}
\begin{enumerate}
  \item Load the Our World in Data vaccination dataset:\\
  \url{https://raw.githubusercontent.com/owid/covid-19-data/master/public/data/vaccinations/vaccinations.csv}
  \item Filter to the 5 countries. Plot \texttt{people\_fully\_vaccinated\_per\_hundred} over time.
\end{enumerate}

\textbf{Panel 4 --- Choropleth map.}
\begin{enumerate}
  \item Using \texttt{geopandas}, create a choropleth of total confirmed cases (or cases per million) across all African countries at the latest available date.
  \item Highlight the 5 focus countries with thicker borders.
\end{enumerate}

\textbf{Layout:}
\begin{minted}{python}
fig, axes = plt.subplots(2, 2, figsize=(18, 14))
# axes[0, 0] -> Panel 1: Cases
# axes[0, 1] -> Panel 2: Deaths
# axes[1, 0] -> Panel 3: Vaccination
# axes[1, 1] -> Panel 4: Choropleth map
\end{minted}

\textbf{Bonus:}
\begin{itemize}
  \item Add a table below the figure showing summary statistics for each country: total cases, total deaths, case fatality rate (\%), vaccination coverage (\%).
  \item Normalize cases and deaths per 100\,000 population to allow fair comparison.
\end{itemize}
\end{exercise}

% ============================================================
\section{Chapter summary}

\begin{itemize}
  \item Health time series have three components: trend, seasonality, and noise.
  \item Rolling averages (\texttt{.rolling()}) smooth daily fluctuations; \texttt{seasonal\_decompose()} separates components formally.
  \item The JHU CSSE dataset provides global COVID-19 case, death, and recovery time series. Use \texttt{.diff()} to convert cumulative counts to daily new cases.
  \item Geospatial data comes in shapefiles (\texttt{.shp}) or GeoJSON (\texttt{.geojson}). \texttt{geopandas} reads both.
  \item Choropleth maps color regions by a health variable --- use \texttt{.plot(column=...)} with an appropriate colormap.
  \item Merging health data with map data requires matching country names --- always check for mismatches.
  \item Small multiples (one map per time period) are more scientifically rigorous than animations.
\end{itemize}

\chapter{Capstone project}
\label{ch:capstone}

\begin{quote}
\emph{``The best way to learn data analysis is to analyze data. Not toy data --- real data, with real messiness, about real health questions.''}
\end{quote}

% ============================================================
\section{Project guidelines}

The capstone project is your opportunity to apply every skill from this course to a real health question. You will work individually or in pairs.

\subsection{Timeline}

\begin{itemize}
  \item \textbf{Week 1.} Choose a project, download the data, perform initial exploration.
  \item \textbf{Week 2.} Complete the analysis, generate figures and tables, draft the report.
  \item \textbf{Week 3.} Finalize the report and prepare a 5-minute oral presentation.
\end{itemize}

\subsection{Expectations}

Your project must include:
\begin{enumerate}
  \item \textbf{A clear health question.} Not ``explore the data'' but ``is diabetes prevalence associated with BMI after adjusting for age and sex?''
  \item \textbf{Real public data.} No simulated datasets. All data sources must be cited.
  \item \textbf{Reproducible code.} A Jupyter notebook that runs from top to bottom without errors.
  \item \textbf{At least 3 visualizations.} Relevant, labeled, and interpretable.
  \item \textbf{At least 1 statistical test or model.} With proper interpretation.
  \item \textbf{A written report.} 5--8 pages following the template below.
  \item \textbf{An oral presentation.} 5 minutes with slides.
\end{enumerate}

\begin{warningbox}
Academic integrity: you may use AI tools (ChatGPT, GitHub Copilot) to help write code, but you must understand every line. During the oral presentation, you will be asked to explain your code and interpret your results. Copy-pasting code you do not understand will be obvious.
\end{warningbox}

% ============================================================
\section{Suggested projects}

Choose one of the five projects below, or propose your own (subject to instructor approval). Each project specifies the objective, dataset(s), required analyses, and expected deliverables.

% ------------------------------------------------------------
\subsection{Project 1: Diabetes risk factor analysis}

\textbf{Objective.} Identify the strongest risk factors for type 2 diabetes and build a predictive model.

\textbf{Datasets.}
\begin{itemize}
  \item \textbf{Pima Indians Diabetes Dataset} (UCI/Kaggle):\\
  \url{https://raw.githubusercontent.com/jbrownlee/Datasets/master/pima-indians-diabetes.data.csv}\\
  768 female patients, 8 features (pregnancies, glucose, blood pressure, skin thickness, insulin, BMI, diabetes pedigree function, age), binary outcome.
  \item \textbf{CDC BRFSS} (Behavioral Risk Factor Surveillance System):\\
  \url{https://www.cdc.gov/brfss/annual_data/annual_data.htm}\\
  Annual survey of 400\,000+ US adults. Download the most recent year's CSV.
\end{itemize}

\textbf{Required analyses.}
\begin{enumerate}
  \item Load and clean the Pima dataset. Handle zeros in glucose, blood pressure, BMI (these are missing values coded as 0).
  \item Exploratory analysis: distributions, correlations, box plots by diabetes status.
  \item Logistic regression: which variables are significant predictors? Report odds ratios with 95\% confidence intervals.
  \item Compare logistic regression with a random forest classifier. Report accuracy, sensitivity, specificity, and AUC-ROC.
  \item (Bonus) Repeat the risk factor analysis on BRFSS data and compare findings.
\end{enumerate}

\textbf{Deliverables.}
\begin{itemize}
  \item Correlation heatmap of all features.
  \item ROC curves comparing logistic regression and random forest.
  \item Table of odds ratios with confidence intervals.
  \item Written interpretation: which risk factors matter most clinically?
\end{itemize}

\begin{clinicalnote}
The Pima Indians dataset is historically important but has ethical considerations. The data was collected from the Akimel O'odham (Pima) community, which has one of the highest rates of type 2 diabetes in the world. When presenting, acknowledge the population and avoid generalizing findings to other groups.
\end{clinicalnote}

% ------------------------------------------------------------
\subsection{Project 2: COVID-19 impact analysis across African countries}

\textbf{Objective.} Analyze how COVID-19 affected African countries differently and explore correlations with health system capacity and socioeconomic factors.

\textbf{Datasets.}
\begin{itemize}
  \item \textbf{JHU CSSE COVID-19 data}:\\
  \url{https://github.com/CSSEGISandData/COVID-19/tree/master/csse_covid_19_data/csse_covid_19_time_series}\\
  Global confirmed cases, deaths, and recovered (daily time series).
  \item \textbf{Our World in Data} (vaccination, testing):\\
  \url{https://raw.githubusercontent.com/owid/covid-19-data/master/public/data/owid-covid-data.csv}\\
  Comprehensive dataset with cases, deaths, vaccinations, testing, hospital data, and demographic indicators.
  \item \textbf{World Bank Open Data} (GDP, health expenditure, hospital beds):\\
  \url{https://data.worldbank.org/}\\
  Use indicator codes: \texttt{SH.XPD.CHEX.PC.CD} (health expenditure per capita), \texttt{SH.MED.BEDS.ZS} (hospital beds per 1\,000), \texttt{NY.GDP.PCAP.CD} (GDP per capita).
  \item \textbf{Africa CDC}:\\
  \url{https://africacdc.org/covid-19/}\\
  Continental-level dashboards and reports.
\end{itemize}

\textbf{Required analyses.}
\begin{enumerate}
  \item Select 10--15 African countries across different subregions (West, East, Southern, North, Central Africa).
  \item Plot epidemic curves (7-day rolling average) for each country. Identify wave patterns.
  \item Compute key metrics: total cases per million, deaths per million, case fatality rate, vaccination rate.
  \item Merge with World Bank indicators. Compute correlations between COVID-19 outcomes and health system capacity (health expenditure, hospital beds, physician density).
  \item Build a choropleth map of case fatality rates across Africa.
  \item (Bonus) Run a multiple regression: does health expenditure per capita predict case fatality rate after controlling for median age and GDP?
\end{enumerate}

\textbf{Deliverables.}
\begin{itemize}
  \item Epidemic curves for all selected countries (small multiples or grouped plot).
  \item Scatter plot of health expenditure vs.\ case fatality rate.
  \item Choropleth map of Africa showing COVID-19 burden.
  \item Summary table with key metrics for all countries.
\end{itemize}

% ------------------------------------------------------------
\subsection{Project 3: Maternal mortality determinants}

\textbf{Objective.} Investigate which health system and socioeconomic factors are associated with maternal mortality across countries.

\textbf{Datasets.}
\begin{itemize}
  \item \textbf{WHO Global Health Observatory} --- Maternal mortality ratio (MMR):\\
  \url{https://apps.who.int/gho/athena/api/GHO/MDG_0000000026?format=csv}
  \item \textbf{WHO GHO} --- Skilled birth attendance:\\
  \url{https://apps.who.int/gho/athena/api/GHO/MDG_0000000025?format=csv}
  \item \textbf{WHO GHO} --- Antenatal care coverage (4+ visits):\\
  \url{https://apps.who.int/gho/athena/api/GHO/WHS4_154?format=csv}
  \item \textbf{World Bank} --- Female literacy rate, GDP per capita, health expenditure:\\
  \url{https://data.worldbank.org/}
  \item \textbf{DHS Program} (Demographic and Health Surveys) --- survey-level indicators:\\
  \url{https://www.statcompiler.com/}\\
  Provides subnational data on maternal health, contraceptive use, and antenatal care for African and Asian countries.
\end{itemize}

\textbf{Required analyses.}
\begin{enumerate}
  \item Load MMR data and filter to the most recent year. Identify the 20 countries with the highest MMR.
  \item Merge with skilled birth attendance, antenatal care coverage, and socioeconomic indicators.
  \item Exploratory analysis: scatter plots of MMR vs.\ each predictor. Compute Pearson and Spearman correlations.
  \item Multiple linear regression: which factors significantly predict MMR? Check model assumptions (residual plots, normality).
  \item Create a choropleth map of MMR across Africa.
  \item (Bonus) Compare MMR trends over time (2000--2020) for 5 countries that implemented specific maternal health programs.
\end{enumerate}

\textbf{Deliverables.}
\begin{itemize}
  \item Scatter plots with regression lines (MMR vs.\ skilled birth attendance, MMR vs.\ health expenditure).
  \item Regression table with coefficients, standard errors, p-values, and $R^2$.
  \item Choropleth map of maternal mortality in Africa.
  \item Discussion of policy implications: which interventions could reduce MMR?
\end{itemize}

\begin{clinicalnote}
Maternal mortality is one of the starkest health inequities in the world. Sub-Saharan Africa accounts for approximately 70\% of global maternal deaths. Most of these deaths are preventable with skilled birth attendants, emergency obstetric care, and adequate antenatal visits.
\end{clinicalnote}

% ------------------------------------------------------------
\subsection{Project 4: Heart disease prediction model comparison}

\textbf{Objective.} Compare multiple machine learning models for predicting heart disease and identify the most important clinical predictors.

\textbf{Dataset.}
\begin{itemize}
  \item \textbf{UCI Heart Disease Dataset} (Cleveland):\\
  \url{https://archive.ics.uci.edu/ml/machine-learning-databases/heart-disease/processed.cleveland.data}\\
  303 patients, 13 clinical features (age, sex, chest pain type, resting BP, cholesterol, fasting blood sugar, resting ECG, max heart rate, exercise-induced angina, ST depression, slope, number of major vessels, thalassemia), binary target (presence of heart disease).
  \item Column names: \texttt{age, sex, cp, trestbps, chol, fbs, restecg, thalach, exang, oldpeak, slope, ca, thal, target}.
\end{itemize}

\textbf{Required analyses.}
\begin{enumerate}
  \item Load and clean the data. Handle missing values (coded as \texttt{?} in the original file).
  \item Exploratory analysis: distributions of each feature by heart disease status.
  \item Train-test split (80/20). Use stratified splitting to maintain class balance.
  \item Train and evaluate 4 models:
  \begin{itemize}
    \item Logistic regression
    \item Decision tree
    \item Random forest
    \item K-nearest neighbors (KNN)
  \end{itemize}
  \item For each model: report accuracy, precision, recall, F1-score, and AUC-ROC.
  \item Plot ROC curves for all 4 models on the same figure.
  \item Feature importance: which clinical variables are the strongest predictors?
\end{enumerate}

\textbf{Deliverables.}
\begin{itemize}
  \item Comparison table of model performance metrics.
  \item ROC curves for all models on one plot.
  \item Feature importance bar chart (from random forest).
  \item Confusion matrix for the best-performing model.
  \item Clinical interpretation: do the top predictors align with cardiology guidelines?
\end{itemize}

\begin{pythontip}
Use \texttt{sklearn.model\_selection.cross\_val\_score} with 5-fold cross-validation instead of a single train-test split. This gives a more reliable estimate of model performance, especially with a small dataset like this one (303 patients).
\end{pythontip}

% ------------------------------------------------------------
\subsection{Project 5: Malaria burden and intervention effectiveness}

\textbf{Objective.} Analyze trends in malaria burden across African countries and evaluate whether key interventions (bed nets, indoor spraying, treatment) are associated with declining incidence.

\textbf{Datasets.}
\begin{itemize}
  \item \textbf{WHO GHO} --- Malaria incidence (per 1\,000 population at risk):\\
  \url{https://apps.who.int/gho/athena/api/GHO/MALARIA_INCIDENCE?format=csv}
  \item \textbf{WHO GHO} --- Malaria mortality rate:\\
  \url{https://apps.who.int/gho/athena/api/GHO/MALARIA_DEATHS_PER100000?format=csv}
  \item \textbf{WHO World Malaria Report data} --- intervention coverage (ITN use, IRS coverage, ACT treatment):\\
  \url{https://www.who.int/teams/global-malaria-programme/reports/world-malaria-report}\\
  Annex data tables are available as Excel downloads.
  \item \textbf{The Malaria Atlas Project}:\\
  \url{https://malariaatlas.org/}\\
  Subnational prevalence estimates and raster maps.
  \item \textbf{DHS Program}:\\
  \url{https://www.statcompiler.com/}\\
  Survey-based data on ITN ownership and use, IRS coverage, and treatment-seeking behavior.
\end{itemize}

\textbf{Required analyses.}
\begin{enumerate}
  \item Load malaria incidence and mortality data. Filter to sub-Saharan African countries.
  \item Plot incidence trends (2000--2022) for 10 countries with the highest burden.
  \item Merge with intervention coverage data (ITN use, IRS, ACT access).
  \item Compute correlations between changes in intervention coverage and changes in incidence over time (i.e., does increasing ITN use predict decreasing incidence?).
  \item Create a choropleth map of current malaria incidence across Africa.
  \item Linear regression: predict change in malaria incidence from change in intervention coverage, controlling for GDP per capita and urbanization.
  \item (Bonus) Create a before-and-after comparison for countries that scaled up ITN distribution after the Global Fund grants.
\end{enumerate}

\textbf{Deliverables.}
\begin{itemize}
  \item Time series plot of malaria incidence for 10 countries (2000--2022).
  \item Scatter plot of ITN use change vs.\ incidence change with regression line.
  \item Choropleth map of malaria incidence across Africa.
  \item Regression table with interpretation.
  \item Policy discussion: which interventions show the strongest association with declining malaria?
\end{itemize}

\begin{clinicalnote}
Between 2000 and 2015, malaria mortality in Africa declined by approximately 44\%, largely driven by the scale-up of ITNs, IRS, and artemisinin-based combination therapies (ACTs). Since 2015, progress has stalled, making this analysis timely and policy-relevant.
\end{clinicalnote}

% ============================================================
\section{Report template}

Your written report should follow this structure (5--8 pages, excluding code appendix):

\subsection{1. Introduction (0.5--1 page)}

\begin{itemize}
  \item Health context: why does this question matter?
  \item Specific research question or objective.
  \item Brief description of the dataset(s) used.
\end{itemize}

\subsection{2. Data description (1 page)}

\begin{itemize}
  \item Source(s) and citation(s).
  \item Number of observations and variables.
  \item Key variables: name, type, unit, and clinical meaning.
  \item Summary statistics table (\texttt{.describe()}).
  \item Missing data: how much, which variables, how handled.
\end{itemize}

\subsection{3. Methods (1 page)}

\begin{itemize}
  \item Data cleaning steps (handling missing values, recoding variables, outlier treatment).
  \item Statistical tests used and why (e.g., ``we used the chi-squared test because both variables are categorical'').
  \item Models used (logistic regression, random forest, etc.) with hyperparameters.
  \item Validation strategy (train-test split, cross-validation).
\end{itemize}

\subsection{4. Results (1.5--2 pages)}

\begin{itemize}
  \item Figures and tables with captions.
  \item Report test statistics, p-values, confidence intervals.
  \item Model performance metrics.
  \item Do not interpret here --- just present the numbers.
\end{itemize}

\subsection{5. Discussion (1--1.5 pages)}

\begin{itemize}
  \item What do the results mean clinically?
  \item How do your findings compare with published literature?
  \item What surprised you?
  \item What are the practical implications?
\end{itemize}

\subsection{6. Limitations (0.5 page)}

\begin{itemize}
  \item Data quality issues (missing data, measurement error, selection bias).
  \item Generalizability: does the dataset represent the population you care about?
  \item What analyses would you do with more time or data?
\end{itemize}

\begin{warningbox}
Every figure must have a title, labeled axes, and a caption explaining what the reader should see. Every table must have a header row and a caption. Figures without labels are the most common reason for lost marks.
\end{warningbox}

% ============================================================
\section{Presentation guidelines}

You will give a 5-minute oral presentation followed by 2--3 minutes of questions.

\subsection{Structure}

\begin{enumerate}
  \item \textbf{Slide 1: Title.} Project title, your name, date.
  \item \textbf{Slide 2: Motivation.} Why does this health question matter? One or two key facts.
  \item \textbf{Slide 3: Data.} What data did you use? How many observations? What are the key variables?
  \item \textbf{Slide 4--5: Key results.} Your 2--3 most important figures or tables. Walk the audience through each one.
  \item \textbf{Slide 6: Conclusions.} One or two take-home messages. What should the audience remember?
\end{enumerate}

\subsection{Tips}

\begin{itemize}
  \item \textbf{5 minutes is short.} Practice with a timer. Cut anything that is not essential.
  \item \textbf{One message per slide.} Do not crowd slides with text.
  \item \textbf{Speak to the audience, not the screen.}
  \item \textbf{Anticipate questions.} Know your limitations and be ready to discuss them.
  \item \textbf{No code on slides.} Show results, not the code that produced them.
\end{itemize}

\begin{pythontip}
Save your figures as high-resolution PNGs for slides:
\begin{minted}{python}
fig.savefig("figure1.png", dpi=300, bbox_inches="tight")
\end{minted}
This avoids blurry screenshots and ensures your plots look professional in a presentation.
\end{pythontip}

% ============================================================
\section{Grading rubric}

The capstone project is graded out of 100 points:

\begin{center}
\begin{tabular}{lcp{8cm}}
\toprule
\textbf{Component} & \textbf{Points} & \textbf{Criteria} \\
\midrule
Research question & 10 & Clear, specific, and health-relevant \\
Data handling & 15 & Proper loading, cleaning, missing value treatment, and documentation \\
Exploratory analysis & 15 & Appropriate summary statistics, distributions, and at least 3 well-labeled visualizations \\
Statistical analysis & 20 & Correct choice and application of tests or models; proper interpretation of p-values, confidence intervals, and effect sizes \\
Report quality & 20 & Follows the template; clear writing; figures have titles, labels, and captions; limitations honestly discussed \\
Code quality & 10 & Reproducible notebook that runs top-to-bottom; comments explain key steps; no unnecessary code \\
Oral presentation & 10 & Clear, within time limit, answers questions confidently \\
\midrule
\textbf{Total} & \textbf{100} & \\
\bottomrule
\end{tabular}
\end{center}

\subsection{Grade boundaries}

\begin{itemize}
  \item \textbf{A (90--100):} Excellent analysis with insightful interpretation. Code is clean and well-documented. Goes beyond the minimum requirements.
  \item \textbf{B (75--89):} Solid analysis with correct methodology. Minor issues in presentation or interpretation.
  \item \textbf{C (60--74):} Meets basic requirements but has notable gaps (e.g., missing visualizations, superficial interpretation, code errors).
  \item \textbf{D/F ($<60$):} Incomplete analysis, major methodological errors, or non-reproducible code.
\end{itemize}

% ============================================================
\section{Getting started: a checklist}

\begin{exercise}
Before you leave today, complete this checklist:
\begin{enumerate}
  \item Choose your project (1--5 or a custom proposal).
  \item Download your primary dataset. Load it into a Jupyter notebook and run \texttt{.head()}, \texttt{.shape}, \texttt{.info()}, and \texttt{.describe()}.
  \item Write your research question in one sentence.
  \item List 3 visualizations you plan to create.
  \item List 1 statistical test or model you plan to use.
  \item Identify potential problems: missing data, messy column names, multiple files to merge.
  \item Create a project folder with the structure:
\begin{minted}{text}
capstone_project/
    data/
        raw/          # original downloaded files
        cleaned/      # processed data
    notebooks/
        01_exploration.ipynb
        02_analysis.ipynb
    figures/
    report.pdf
\end{minted}
  \item Submit your project choice and research question to the instructor by the end of the week.
\end{enumerate}
\end{exercise}

% ============================================================
\section{Chapter summary}

\begin{itemize}
  \item The capstone project integrates all course skills: data loading, cleaning, exploration, visualization, statistical analysis, and interpretation.
  \item Five suggested projects cover diabetes, COVID-19, maternal mortality, heart disease, and malaria --- all using real, publicly available datasets.
  \item Each project requires a clear health question, reproducible code, at least 3 visualizations, and at least 1 statistical test or model.
  \item The report follows a standard structure: Introduction, Data, Methods, Results, Discussion, Limitations.
  \item Presentations are 5 minutes --- focus on motivation, key results, and conclusions.
  \item The grading rubric rewards clear thinking and honest interpretation over complex methodology.
\end{itemize}


% ============================================================
\backmatter

\chapter*{Appendix A: Python installation guide}
\addcontentsline{toc}{chapter}{Appendix A: Python installation guide}

\section*{Option 1: Google Colab (recommended for beginners)}
Go to \url{https://colab.research.google.com}. Sign in with a Google account. No installation needed. All libraries used in this course are pre-installed.

\section*{Option 2: Anaconda (local installation)}
Download Anaconda from \url{https://www.anaconda.com/download}. Install with default settings. Open Jupyter Notebook from Anaconda Navigator.

\section*{Required libraries}
\begin{minted}{bash}
pip install pandas matplotlib seaborn scipy scikit-learn lifelines geopandas
\end{minted}

\chapter*{Appendix B: Health data sources}
\addcontentsline{toc}{chapter}{Appendix B: Health data sources}

\begin{longtable}{p{4cm}p{8cm}p{3cm}}
\toprule
\textbf{Source} & \textbf{Description} & \textbf{URL} \\
\midrule
WHO GHO & Global health indicators (190+ countries, 1000+ indicators) & \url{who.int/data/gho} \\
Gapminder & Health, income, population by country and year & \url{gapminder.org/data} \\
Our World in Data & COVID-19, mortality, disease burden, vaccination & \url{ourworldindata.org} \\
CDC NHANES & US national health survey (labs, questionnaires, exams) & \url{cdc.gov/nchs/nhanes} \\
UCI ML Repository & Curated ML datasets (heart disease, diabetes, cancer) & \url{archive.ics.uci.edu} \\
World Bank & Health expenditure, life expectancy, maternal mortality & \url{data.worldbank.org} \\
Africa CDC & African disease surveillance and outbreak data & \url{africacdc.org} \\
JHU COVID-19 & Global COVID-19 cases, deaths, recoveries (time series) & \url{github.com/CSSEGISandData} \\
DHS Program & Demographic and Health Surveys (50+ African countries) & \url{dhsprogram.com} \\
\bottomrule
\end{longtable}

\bibliographystyle{plainnat}
\bibliography{references}

\printindex

\end{document}
